% !TeX root = ../../Book2.tex
% 纲要标题：## 第12章　链条件、长度与直和分解
% 纲要位置：工作记录/计划纲要.md:L1434

\chapter{链条件、长度与直和分解}
\label{b2:ch:12}
\BookChapterNumber{b2:ch:12}

\begin{chapterabstract}
链条件把有限维线性代数中“不能无限增加或减少维数”的现象推广到一般模。本章证明 Noether 条件与子模有限生成的等价性，说明升降链条件在短正合列中的传递及非交换环上的左右区别。合成列与模的 Jordan–Hölder 定理给出长度及其可加性；Fitting 引理把有限长度模分解为自同态作用的幂零分量与可逆分量，并由此证明不可分解模的自同态环局部。最后证明 Krull–Schmidt 分解的存在唯一性，用具体反例说明有限生成或单独的 Noether 条件为何不足。
\end{chapterabstract}

\begin{learninggoals}
\begin{itemize}
\item 区分有限生成、升链条件、降链条件与有限长度，能通过子模、商模及短正合列判断这些性质。
\item 证明模的 Jordan–Hölder 定理，计算长度和简单因子的重数，并区别长度与基域维数。
\item 证明 Fitting 分解，使用局部环判据辨认不可分解模，说明幂等元和投影细分的关系。
\item 在 Krull–Schmidt 唯一性证明中从恒等映射的有限和中找到可逆项，比较分量并交换补模，理解结论的适用条件及失效例子。
\end{itemize}
\end{learninggoals}

整数环 \(\mathbb Z\) 的每个理想都由一个整数生成，因而子模的生成问题十分简单；但
\[
\mathbb Z\supsetneq2\mathbb Z\supsetneq4\mathbb Z\supsetneq8\mathbb Z\supsetneq\cdots
\]
永不停止。它说明升链条件与降链条件不能混为一谈。有限维向量空间同时满足两者，是因为严格变化每次都改变维数；一般模需要分别研究这两种有限性。

\ProseParagraphBreak

当升降过程都能停止时，一个模可以被逐层剖开，记录简单的商层；但仅知道这些层是什么，并不能自动把它们重新拼成直和。Jordan–Hölder 定理确定这些商层的同构类型和重数，Krull–Schmidt 定理确定实际的不可分解直和分量，两者回答不同的问题。为证明后一种唯一性，还需要用链条件控制自同态反复作用后的核与像，由 Fitting 分解得到不可分解模的局部自同态环。

\ProseParagraphBreak

本章需要\chapref{b2:ch:02}的子模与同构定理、\chapref{b2:ch:03}的正合列与直和，以及\chapref{b2:ch:11}关于幂等自同态与模分解的结果。个别算例还会用到\chapref{b2:ch:04}的主理想整环模分类和\chapref{b2:ch:07}的投射模判据。局部环在这里由非单位理想定义。

\ProseParagraphBreak
交换环上的 Noether 模判据可参阅 Roman 的《Advanced Linear Algebra》\cite[Chapter 5, Theorem 5.7, pp. 133--134]{Roman2008AdvancedLinearAlgebra}。有限维表示的合成列与不可分解分解，可参阅 Etingof 等人的《Introduction to Representation Theory》\cite[Theorem 3.7.1, pp. 50--51; Theorem 3.8.1, pp. 51--52]{EtingofEtAl2011RepresentationTheory}。本章讨论一般环上的有限长度模，不预设代数闭底域。

% !TeX root = ../../Book2.tex
% 纲要标题：### 12.1 链条件
% 纲要位置：工作记录/计划纲要.md:L1436

\section{链条件}
\label{b2:sec:1201}

有限维向量空间的严格子空间链必在有限步后终止，因为每次严格包含都改变维数。一般模没有这样的维数，但仍可直接要求子模链终止。本节区分升链与降链两种条件，并说明它们怎样通过子模、商模和短正合列传递。除另外说明外，\(R\) 为含单位元的环，模均为幺性左 \(R\) 模，不假设环交换。

% 纲要要点（供目录核对）：
%
% * Noether 模与 Artin 模
% * 子模、商模与短正合列中的传递
% * 环的左右链条件
%

\subsection{Noether 模与 Artin 模}
\label{b2:subsec:1201:01}

\begin{definition}\label{b2:def:noetherian-artinian-module}
设 \(R\) 为含幺环、\(M\) 为幺性左 \(R\) 模。若 \(M\) 的每条子模升链
\[
N_1\subseteq N_2\subseteq\cdots
\]
都最终稳定，即存在 \(r\) 使 \(N_i=N_r\) 对所有 \(i\ge r\) 成立，则称 \(M\) 满足\term[sheng-lian-tiaojian]{升链条件}[ascending chain condition]，也称 \(M\) 为 \term[Noether-mo]{Noether 模}[Noetherian module]。若每条子模降链 \(N_1\supseteq N_2\supseteq\cdots\) 都最终稳定，则称 \(M\) 满足\term[jiang-lian-tiaojian]{降链条件}[descending chain condition]，也称 \(M\) 为 \term[Artin-mo]{Artin 模}[Artinian module]。两种条件分别简称 ACC 与 DCC。
\end{definition}

“最终稳定”不是说所有链都有一个共同的长度上界。例如 \(\mathbb Z\) 的每个子模为主理想，升链最终稳定；但
\[
0\subsetneq 2^n\mathbb Z\subsetneq2^{n-1}\mathbb Z
\subsetneq\cdots\subsetneq\mathbb Z
\]
可以任意长。这个模也不满足降链条件，因为 \(\mathbb Z\supsetneq2\mathbb Z\supsetneq4\mathbb Z\supsetneq\cdots\) 永不稳定。

\begin{proposition}\label{b2:prop:chain-maximal-minimal}
设 \(R\) 为含幺环、\(M\) 为幺性左 \(R\) 模。\(M\) 满足升链条件，当且仅当每个非空子模集合在包含关系下有极大元。把升链、极大元分别换成降链、极小元，也有同样的等价。
\end{proposition}
\begin{proof}
若某个非空子模集合没有极大元，从其中任取 \(N_1\)，再逐次取严格包含前一项的 \(N_{i+1}\)，得到不稳定的升链，矛盾。反之，升链的全部项组成非空集合，若其中某项 \(N_r\) 极大，则以后各项包含它，只能都等于它。降链的论证将包含方向反转即可：无极小元时逐次严格缩小；有极小项时后续项与之相等。这里的逐次选取沿用本书的选择公理约定。
\end{proof}

这里的极大元只要求不能在给定子模集合中继续严格增大，不要求包含这个集合中的所有成员。例如，域 \(k\) 上 \(k^2\) 中两个不同的一维子空间，在由它们组成的集合中都为极大元，却都不是最大元。在一条链中，各项可比，极大项才自动成为最大项；极小与最小的区别同样如此。

\begin{proposition}\label{b2:prop:noetherian-submodules-fg}
设 \(R\) 为含幺环、\(M\) 为幺性左 \(R\) 模。\(M\) 是 Noether 模，当且仅当它的每个子模都作为左 \(R\) 模有限生成。
\end{proposition}
\begin{proof}
若某个子模 \(N\) 不有限生成，就可以不断加入旧生成元无法表示的新元素。先取 \(x_1\in N\)，再从 \(N\setminus(Rx_1+\cdots+Rx_i)\) 取 \(x_{i+1}\)，得到严格升链
\[
Rx_1\subsetneq Rx_1+Rx_2\subsetneq\cdots.
\]
这与升链条件矛盾，故每个子模有限生成。零子模由空集生成，不需另选元素。

反过来，给定升链 \((N_i)\)，其并 \(N=\bigcup_iN_i\) 为子模：两个元素同属某个较大的 \(N_i\)，和与标量倍都仍在其中。取 \(N\) 的有限生成元 \(x_1,\ldots,x_s\)。每个生成元在链的某一项中，有限多个这样的项有共同的后继项 \(N_r\)，所以全部生成元已经同时落在 \(N_r\) 中；零子模的情形可任取 \(r\)。于是 \(N\subseteq N_r\subseteq N\)，后续各项只能等于 \(N_r\)，所以升链稳定。
\end{proof}

特别地，Noether 模本身有限生成，反向却依赖于环。例如 \(R=k[X_1,X_2,\ldots]\) 是循环左 \(R\) 模，但理想链
\((X_1)\subsetneq(X_1,X_2)\subsetneq\cdots\) 严格上升。严格性可令前 \(i\) 个不定元取零来核验：\(X_{i+1}\) 不可能属于前面的理想。有限生成只控制整个模的生成元，并不自动控制全部子模。

\ProseParagraphBreak
Artin 条件与升链条件的差别可在交换群中直接看到。固定素数 \(p\)，令
\[
C_{p^\infty}=\mathbb Z[1/p]/\mathbb Z,\qquad
C_{p^n}=\langle p^{-n}+\mathbb Z\rangle.
\]
每个元素的阶为某个 \(p\) 的幂。阶为 \(p^m\) 的元素可写成 \(a/p^m+\mathbb Z\)，其中分子与 \(p\) 互素，因而它生成 \(C_{p^m}\)。被 \(p^n\) 消去的全部元素恰好组成 \(C_{p^n}\)：这正要求约分后的分母不超过 \(p^n\)。所以，若子群 \(H\) 中的元素阶无界，对每个 \(n\) 都能找到阶至少为 \(p^n\) 的元素，其生成的循环群包含 \(C_{p^n}\)，于是 \(H=C_{p^\infty}\)。若非零子群 \(H\) 的元素阶有界，取其中达到最大阶 \(p^m\) 的元素，便有 \(C_{p^m}\subseteq H\)；同时 \(H\) 的全部元素都被 \(p^m\) 消去，所以 \(H\subseteq C_{p^m}\)。这就证明每个真子群都是某个有限循环群 \(C_{p^m}\) 或零群。

\ProseParagraphBreak
因此，\(C_{p^\infty}\) 中的降链一旦离开全群，便落在一个有限群中，此后必稳定，所以它是 Artin \(\mathbb Z\) 模。另一方面，\(C_p\subsetneq C_{p^2}\subsetneq\cdots\) 永不稳定，它不是 Noether 模；任何有限组元素又同时落在某个 \(C_{p^n}\) 中，不能生成全群，所以它也不是有限生成模。与上面 \(\mathbb Z\) 为 Noether 模而非 Artin 模的例子合起来可见，这两种链条件互不蕴含，Artin 条件也不保证有限生成。

\subsection{链条件在短正合列中的传递}
\label{b2:subsec:1201:02}

在子模和商模上同时检测一条链，可以把中间模的增长或缩小分成两部分。对于可比的子模 \(L\subseteq L'\)，若增长既没有在固定子模 \(A\) 内留下新部分，也没有在商模中留下新像，那么实际没有增长。下面的引理把这个比较写成精确的交与像条件。

\begin{lemma}\label{b2:lem:submodule-intersection-image}
设 \(R\) 为含幺环、\(B\) 为幺性左 \(R\) 模、\(A\subseteq B\) 为子模，\(q:B\rightarrow B/A\) 为商映射。若子模 \(L\subseteq L'\subseteq B\) 满足
\[
L\cap A=L'\cap A,\qquad q(L)=q(L'),
\]
则 \(L=L'\)。
\end{lemma}
\begin{proof}
给定 \(x\in L'\)，由像相同可取 \(y\in L\) 使 \(q(x)=q(y)\)。于是 \(x-y\in L'\cap A=L\cap A\subseteq L\)，所以 \(x=y+(x-y)\in L\)。这证明反向包含，得到相等。
\end{proof}

\begin{theorem}\label{b2:thm:chain-conditions-exact}
设 \(R\) 为含幺环，\(0\rightarrow A\rightarrow B\rightarrow C\rightarrow0\) 为幺性左 \(R\) 模的短正合列，则 \(B\) 为 Noether 模，当且仅当 \(A,C\) 均为 Noether 模；对 Artin 模也有相同结论。
\end{theorem}
\begin{proof}
将 \(A\) 识别为 \(B\) 的子模，\(C\) 识别为 \(B/A\)。若 \(B\) 的子模升链稳定，\(A\) 中的升链也是 \(B\) 中的升链，故稳定；商模 \(C\) 中的升链拉回为包含 \(A\) 的升链，由\cref{b2:thm:module-correspondence}的子模对应，其稳定性也成立。

反过来，给定 \(B_1\subseteq B_2\subseteq\cdots\)，同时观察它在 \(A\) 内的交链 \(B_i\cap A\) 与在 \(C\) 中的像链 \(q(B_i)\)。若 \(A,C\) 均为 Noether 模，取两条链稳定起点中的较大者 \(r\)。对每个 \(i\ge r\)，\(B_r\subseteq B_i\) 有相同交与像，\cref{b2:lem:submodule-intersection-image}给出 \(B_r=B_i\)，所以两端停止变化后，中间链也停止变化。

对降链，子模继承和商模拉回仍然有效。若两端满足降链条件，取交链和像链的共同稳定起点 \(r\)，对 \(i\ge r\) 应用该引理于 \(B_i\subseteq B_r\)，又得相等。因此 Artin 情形的两方向都成立。
\end{proof}

由 \(0\rightarrow M\rightarrow M\oplus N\rightarrow N\rightarrow0\)，有限直和满足某一种链条件，当且仅当每个分量都满足。这一结论不能推广为无限直和：非零模 \(M_i\) 的可数直和有严格升链 \(\bigoplus_{i\le n}M_i\)，也有严格降链 \(\bigoplus_{i\ge n}M_i\)。两种严格性都可用某个非零坐标元素核验。

\ProseParagraphBreak
引理中的“有相同交与像”必须配合包含关系使用。两个不同的补模可以与 \(A\) 都交于零、都满射到 \(B/A\)，却仍然不同。例如在 \(k^2\) 中，以 \(A=ke_1\) 为固定子空间，\(ke_2\) 与 \(k(e_1+e_2)\) 就有相同的交与像。链条件证明中，各项彼此可比，才保证引理适用。

\subsection{环的左右链条件}
\label{b2:subsec:1201:03}

\begin{definition}\label{b2:def:left-right-chain-ring}
设 \(R\) 为含幺环。若左正则模 \({}_RR\) 是 Noether 模，则称 \(R\) 为\term[zuo-Noether-huan]{左 Noether 环}[left Noetherian ring]；若 \({}_RR\) 是 Artin 模，则称 \(R\) 为\term[zuo-Artin-huan]{左 Artin 环}[left Artinian ring]。相应地，对右正则模 \(R_R\) 定义右 Noether 环和右 Artin 环。它们分别要求左理想链或右理想链满足相应条件。
\end{definition}

在交换环中，左右理想相同，通常省略方向。非交换环中则必须保留它。\cref{b2:prop:noetherian-submodules-fg}说明，左 Noether 环等价于每个左理想都作为左模有限生成；这没有自动给出右理想的有限生成性。

\begin{proposition}\label{b2:prop:ring-chain-finitely-generated}
设 \(R\) 为含幺环。\(R\) 是左 Noether 环，当且仅当每个有限生成幺性左 \(R\) 模都为 Noether 模。把 Noether 换成 Artin，同样成立。
\end{proposition}
\begin{proof}
若左正则模满足所讨论的链条件，\cref{b2:thm:chain-conditions-exact}说明有限直和 \(R^n\) 也满足。任意有限生成模都有满同态 \(R^n\rightarrow M\)，商模继承链条件，所以 \(M\) 满足。反向取有限生成左模 \(M=R\) 即得环本身的条件。两个论证对升链和降链分别有效。
\end{proof}

要让左右链条件不同，可以让同一块的左侧使用一个大标量域，右侧却只使用它的子域。令 \(k=\mathbb Q\)、\(K=\mathbb Q(t)\)，取三角矩阵环
\[
R=\left\{\,
\begin{pmatrix}a&u\\0&b\end{pmatrix}
\ \middle|\ a,u\in K,\ b\in k
\,\right\}.
\]
记 \(e_{11}=\begin{psmallmatrix}1&0\\0&0\end{psmallmatrix}\)、\(e_{22}=\begin{psmallmatrix}0&0\\0&1\end{psmallmatrix}\)，并令
\[
J=\left\{\,
\begin{pmatrix}0&u\\0&0\end{pmatrix}\ \middle|\ u\in K
\,\right\},
\]
将其中的元素记为 \(j(u)\)。对 \(r=\begin{psmallmatrix}a&v\\0&b\end{psmallmatrix}\in R\)，直接相乘得到
\[
rj(u)=j(au),\qquad j(u)r=j(ub).
\]
因此 \(J\) 左侧可以使用全部 \(K\) 标量，右侧却只能使用 \(k\) 标量。左模 \(Re_{11}\) 由左上位置的 \(K\) 组成，\(J\subseteq Re_{22}\) 也以 \(K\) 中的标量作左作用；这两个左模的任意非零元素都生成全模，所以没有非零真子模。商模 \(Re_{22}/J\cong k\) 只留下右下坐标，其左作用为乘 \(b\in k\)，也没有非零真子模。这三个模都满足两种链条件，故由短正合列传递和按列分解 \(R=Re_{11}\oplus Re_{22}\)，\(R\) 同时为左 Noether 环和左 Artin 环。

\ProseParagraphBreak
在右侧，上面的公式还说明 \(J_R\) 的子模恰好对应 \(K\) 的 \(k\) 线性子空间：任意 \(b\in k\) 都能作为右下坐标出现，所以右作用下封闭正是对这些标量倍封闭。取
\[
U_n=\operatorname{span}_k\{1,t,\ldots,t^n\},\qquad
W_n=\operatorname{span}_k\{t^j:j\ge n\},
\]
得到严格上升的 \(U_n\) 与严格下降的 \(W_n\)，因为不定元的不同幂在 \(k\) 上线性无关。由 \(j\) 将它们送入 \(J\)，便得到两条真正的右理想链，而不只是向量子空间链。因此这个 \(R\) 既不是右 Noether 环，也不是右 Artin 环；即使两种左链条件同时成立，也不能略去左右方向。若将 \(K/k\) 换成有限扩张，两侧都能具有有限的简单商层，但层数仍可能不同，\cref{b2:prop:triangular-left-right-length}将给出具体计算。

% !TeX root = ../../Book2.tex
% 纲要标题：### 12.2 有限长度
% 纲要位置：工作记录/计划纲要.md:L1442

\section{有限长度}
\label{b2:sec:1202}

升链条件阻止不断增加新子模，降链条件阻止不断缩小旧子模；同时具有这两种条件时，模可以在有限步内逐层取商，直到各商层都没有非零真子模。本节首先证明这些层的同构类型和重数不依赖于合成列的选择，再据此定义长度。长度计算的是简单商层，原模中相邻商层之间的扩张可以不分裂。

% 纲要要点（供目录核对）：
%
% * 合成列与 Jordan–Hölder 定理
% * 长度的可加性
% * 有限维代数的有限生成模与左右长度
%

\subsection{合成列与 Jordan–Hölder 定理}
\label{b2:subsec:1202:01}

\begin{definition}\label{b2:def:composition-series}
设 \(R\) 为含幺环。若非零幺性左 \(R\) 模 \(S\) 只有 \(0,S\) 两个子模，则称 \(S\) 为\term[jiandan-mo]{简单模}[simple module]。对幺性左 \(R\) 模 \(M\)，取有限子模链
\[
0=M_0\subsetneq M_1\subsetneq\cdots\subsetneq M_n=M,
\]
若其中每个商 \(M_i/M_{i-1}\) 都简单，则称这条链为 \(M\) 的\term[hecheng-lie]{合成列}[composition series]，这些商称为该列的\term[hecheng-yinzi]{合成因子}[composition factor]。若 \(M\) 存在合成列，则称它为\term[youxian-changdu-mo]{有限长度模}[module of finite length]；零模采用没有非零商层的空合成列。
\end{definition}

固定素数 \(p\)。例如 \(\mathbb Z/p\mathbb Z\) 为简单 \(\mathbb Z\) 模；而 \(\mathbb Z/p^2\mathbb Z\) 不简单，因为其中有非零真子模 \(p\mathbb Z/p^2\mathbb Z\)。不过
\[
0\subsetneq p\mathbb Z/p^2\mathbb Z\subsetneq\mathbb Z/p^2\mathbb Z
\]
是一条合成列，两层都同构于 \(\mathbb Z/p\mathbb Z\)。\((\mathbb Z/p\mathbb Z)^2\) 有同样的两个合成因子，却不与前者同构：前者有阶为 \(p^2\) 的元素，后者每个元素都被 \(p\) 零化。用\cref{b2:thm:split-exact-sequence}的语言来看，第一条列对应的短正合列
\[
0\longrightarrow\mathbb Z/p\mathbb Z
\xrightarrow{\ \bar a\mapsto pa+p^2\mathbb Z\ }
\mathbb Z/p^2\mathbb Z
\longrightarrow\mathbb Z/p\mathbb Z\longrightarrow0
\]
不分裂：中间模里被 \(p\) 消去的元素都在 \(p\mathbb Z/p^2\mathbb Z\) 中，取商后为零，不可能给出截面。而以 \((\mathbb Z/p\mathbb Z)^2\) 为中间模、把第一坐标作为子模的相应短正合列分裂，第二坐标就是补模。合成因子记录简单商层，不记录相邻层之间的扩张是否分裂，也不表示这些层都已成为原模的直和分量。

\begin{lemma}\label{b2:lem:composition-series-subquotient}
设 \(R\) 为含幺环、\(M\) 为有限长度幺性左 \(R\) 模。\(M\) 的每个子模和商模都有合成列。若 \(M\) 有一条包含 \(n\) 个简单商层的合成列，则所构造的子模与商模合成列各至多有 \(n\) 层。
\end{lemma}
\begin{proof}
给定合成列 \(0=M_0\subsetneq\cdots\subsetneq M_n=M\) 及子模 \(N\subseteq M\)。把各项与 \(N\) 相交，得到
\[
0=N\cap M_0\subseteq N\cap M_1\subseteq\cdots\subseteq N\cap M_n=N.
\]
包含映射诱导
\[
(N\cap M_i)/(N\cap M_{i-1})\longrightarrow M_i/M_{i-1}.
\]
它单射，因为映到零正好要求代表元落在 \(M_{i-1}\)。像为简单模的子模，所以或为零、或为全体。删去重复项，得到 \(N\) 的合成列。

再把各项映到 \(M/N\)，得到
\[
0=(M_0+N)/N\subseteq(M_1+N)/N\subseteq\cdots\subseteq(M_n+N)/N=M/N.
\]
相邻商是 \(M_i/M_{i-1}\) 的商：映射把 \(m+M_{i-1}\) 送到 \(m+(M_{i-1}+N)\)，在相应商层上满射。因此它或为零、或与原简单模同构；为零正是相邻项重复的情形。删去重复项，即得 \(M/N\) 的合成列。两种过程都不会增加原来的 \(n\) 个商层。
\end{proof}

比较两条合成列时，若它们紧邻全模的子模不同，可以先取这两个子模的交，再把两种最后的简单商交叉配对。这样剩余部分可以归纳比较，末端只交换两个简单因子的次序。

\begin{theorem}[Jordan–Hölder]\label{b2:thm:module-jordan-holder}
设 \(R\) 为含幺环、\(M\) 为有限长度幺性左 \(R\) 模。\(M\) 的任意两条合成列具有相同的层数，并且其合成因子在计入重数后，除排列次序外逐个同构。
\end{theorem}
\begin{proof}
对第一条合成列的层数 \(n\) 作强归纳。\(n=0\) 时 \(M=0\)，只有空列。设 \(n>0\)，记两条合成列中紧邻 \(M\) 的真子模为 \(A,B\)，所以 \(M/A,M/B\) 简单。

若 \(A=B\)，对 \(A\) 的两条合成列应用归纳假设，再在两侧加上相同的商 \(M/A\)，即得结论。现设 \(A\ne B\)。它们都是极大真子模，因此 \(A+B=M\)；否则 \(A+B\) 仍真且包含 \(A\)，会等于 \(A\)，从而 \(B\subseteq A\)，再由 \(B\) 极大得到 \(B=A\)，矛盾。

令 \(C=A\cap B\)。映射 \(A\rightarrow M/B\)、\(a\mapsto a+B\) 的核为 \(C\)，像由 \(A+B=M\) 为全商，因此\cref{b2:thm:module-first-isomorphism}给出 \(A/C\cong M/B\)。同样有 \(B/C\cong M/A\)。这两个商都简单。

由\cref{b2:lem:composition-series-subquotient}，\(C\subseteq A\) 有某条有限合成列，设这条列有 \(r\) 层。将其后接简单商 \(A/C\)，得到 \(A\) 的一条 \(r+1\) 层合成列。第一条原列在 \(A\) 中已有 \(n-1\) 层，故以原列的这一部分应用归纳假设，得到 \(r+1=n-1\)，且两种因子列表相同。于是所取 \(C\) 的合成列恰有 \(n-2\) 层。

把同一条 \(C\) 的合成列后接 \(B/C\)，得到 \(B\) 的一条 \(n-1\) 层合成列。以这条新列作为强归纳中的第一条列，就可与第二条原列在 \(B\) 中的部分比较。第一条 \(M\) 合成列的因子，因而由 \(C\) 的这些因子再加 \(A/C\cong M/B\) 与 \(M/A\) 组成；第二条则由同一组 \(C\) 因子再加 \(B/C\cong M/A\) 与 \(M/B\) 组成。它们只是最后两种因子的次序互换，故层数和全部重数相同。
\end{proof}

对含幺环 \(R\) 上的有限长度幺性左模 \(M\)，模的\term[mo-Jordan-Holder-dingli]{Jordan–Hölder 定理}[Jordan–Hölder theorem for modules]使以下记号有意义：
\[
\ell_R(M)=\text{\(M\) 的任意合成列的层数}.
\]
这个数称为 \(M\) 的\term[mo-changdu]{长度}[length of a module]；环已明确时简写 \(\ell(M)\)。它为零当且仅当 \(M=0\)。对于简单左 \(R\) 模 \(S\)，还可把它在合成列中出现的次数记为 \([M:S]\)，称为相应合成重数；只比较同构类型，不把具体嵌入方式计入这个数。
\symidx{\(\ell_R(M)\)：左 \(R\) 模的长度}
\symidx{\([M:S]\)：简单模 \(S\) 在 \(M\) 中的合成重数}

\subsection{长度的可加性}
\label{b2:subsec:1202:02}

\begin{theorem}\label{b2:thm:finite-length-chain-conditions}
设 \(R\) 为含幺环、\(M\) 为幺性左 \(R\) 模。\(M\) 有限长度，当且仅当它同时为 Noether 模和 Artin 模。
\end{theorem}
\begin{proof}
简单模没有可继续严格变化的子模链，故同时满足两种链条件。若 \(M\) 有合成列，逐层应用\cref{b2:thm:chain-conditions-exact}，便知全部层的扩张 \(M\) 同时满足两种链条件。

反过来，设 \(M\) 同时满足两种条件。升链条件用来保证每一步能找到简单商，降链条件则用来保证这个过程在有限步后结束。若 \(M\ne0\)，由\cref{b2:prop:chain-maximal-minimal}，全部真子模中有极大元 \(M^{(1)}\)，其商简单；若 \(M^{(1)}\ne0\)，再在它内部选取极大真子模。子模继承升链条件，因此每一步都可进行。若永不到零，就得到无限严格降链，违反 Artin 条件。所以经过有限步到达零，反向排列即为合成列。
\end{proof}

\begin{proposition}[长度的可加性]\label{b2:prop:length-additive}
设 \(R\) 为含幺环，\(0\rightarrow A\rightarrow B\rightarrow C\rightarrow0\) 为幺性左 \(R\) 模的短正合列。\(B\) 有限长度当且仅当 \(A,C\) 都有限长度。此时
\[
\ell(B)=\ell(A)+\ell(C),\qquad
[B:S]=[A:S]+[C:S]
\]
对每个简单幺性左 \(R\) 模 \(S\) 都成立，其中 \(\ell\) 表示合成列长度，\([X:S]\) 表示 \(S\) 在 \(X\) 的合成列中出现的重数。
\end{proposition}
\begin{proof}
若 \(B\) 有限长度，子模 \(A\) 与商模 \(C\) 由\cref{b2:lem:composition-series-subquotient}也有合成列。反过来，给定 \(A,C\) 的合成列，把 \(C\) 中的各项通过商映射拉回到 \(B\)，得到从 \(A\) 到 \(B\) 的链，其相邻商与 \(C\) 的各商同构；再在前面接上 \(A\) 的合成列，就得到 \(B\) 的合成列。新列恰好把两端的全部因子合并，Jordan–Hölder 定理保证所计算的长度与重数不依赖这些选择，于是两个公式成立。
\end{proof}

因此有限长度模的真子模和真商模长度都严格变小。对严格链 \(0=L_0\subsetneq L_1\subsetneq\cdots\subsetneq L_r=M\)，各 \(L_i/L_{i-1}\) 非零，长度至少为一，反复使用可加性得到 \(r\le\ell(M)\)。合成列达到这个上界，所以长度也是严格子模链可达到的最大层数。\secref{b2:sec:1201}中的 \(\mathbb Z\) 虽然满足升链条件，却有任意长的有限严格链；有限长度在此提供了真正统一的上界。

\begin{proposition}\label{b2:prop:finite-length-endomorphism}
设 \(R\) 为含幺环、\(M\) 为有限长度幺性左 \(R\) 模。对自同态 \(f:M\rightarrow M\)，单射、满射、可逆这三个条件等价。更一般地，
\[
\ell(M)=\ell(\ker f)+\ell(\operatorname{im}f)
\]
对任何到幺性左 \(R\) 模 \(N\) 的同态 \(f:M\to N\) 成立。
\end{proposition}
\begin{proof}
短正合列 \(0\rightarrow\ker f\rightarrow M\rightarrow\operatorname{im}f\rightarrow0\) 与长度可加性给出公式。若 \(f\) 是自同态且满射，像长度等于 \(\ell(M)\)，故核长度为零，核为零；若它单射，像长度等于 \(\ell(M)\)，再对像的包含取商，商长度为零，故像为全模。模的双射同态有线性逆，因而可逆。可逆显然同时给出单射和满射。
\end{proof}

这个长度公式对应线性代数中维数等于核维数与像维数之和；有限长度自同态的单射与满射等价，也对应有限维线性空间中的结论。例如，对素数 \(p\) 和 \(n\ge1\)，\(M=\mathbb Z/p^n\mathbb Z\) 的链 \(0\subsetneq p^{n-1}M\subsetneq\cdots\subsetneq M\) 有 \(n\) 个简单商层，所以长度为 \(n\)。乘 \(p\) 的核长度为一、像长度为 \(n-1\)，可加性给出 \(n=1+(n-1)\)。在 \(n>1\) 时，这个算子既不单射，也不满射。

\subsection{有限维代数的有限生成模}
\label{b2:subsec:1202:03}

域 \(k\) 上的简单模恰好是一维向量空间，所以对有限维 \(k\) 向量空间 \(V\)，有 \(\ell_k(V)=\dim_kV\)。一般代数上的简单模却可能有更大的 \(k\) 维数；要比较长度和维数，需要同时计算每个简单商层的维数。

\begin{proposition}\label{b2:prop:finite-dimensional-algebra-modules}
设 \(k\) 为域、\(A\) 为有限维结合含幺 \(k\) 代数、\(M\) 为幺性左 \(A\) 模。以下三个条件等价：\(M\) 有限生成；\(M\) 作为 \(k\) 向量空间有限维；\(M\) 作为 \(A\) 模有限长度。若 \(S_1,\ldots,S_r\) 为一条合成列的因子，则
\[
\dim_kM=\sum_{i=1}^r\dim_kS_i,\qquad \ell_A(M)\le\dim_kM.
\]
\end{proposition}
\begin{proof}
有限生成时存在满同态 \(A^n\rightarrow M\)，因 \(A\) 有限维，目标的 \(k\) 维数有限。有限维时，严格子模链也为严格子空间链，维数只能有限次增加或减少，所以 \(M\) 同时满足两种链条件，由\cref{b2:thm:finite-length-chain-conditions}得有限长度。有限长度又蕴含 Noether 条件，故由\cref{b2:prop:noetherian-submodules-fg}得 \(M\) 有限生成，完成等价。

每个简单因子 \(S_i\) 由任一非零元素生成，故是 \(A\) 的商模，具有有限且正的 \(k\) 维数。合成列逐层使用向量空间短正合列的维数公式，给出所列和式；每项至少为一，所以得到长度的不等式。
\end{proof}

长度计算简单商层的个数，维数则是这些商层的 \(k\) 维数之和，所以二者一般不同。取 \(A=M_n(k)\)、\(n\ge1\)，其标准列模 \(k^n\) 简单：任取非零向量，利用某个非零坐标和矩阵单位可以得到每个标准基向量，故它生成全模。因此这个模的 \(A\) 长度为一，而 \(k\) 维数为 \(n\)。对于交换环 \(k[t]\) 上的模 \(k[t]/(p^m)\)，其中 \(m\ge1\)、\(p\in k[t]\) 首一不可约，逐次乘 \(p\) 给出 \(m\) 个同构于 \(k[t]/(p)\) 的简单层；其长度为 \(m\)，\(k\) 维数为 \(m\deg p\)。

\ProseParagraphBreak
代数有限维是上述等价的实质条件。\(k[t]\) 作为自身的模由一生成，却有不稳定降链 \((t)\supsetneq(t^2)\supsetneq\cdots\)，所以不有限长度；域 \(k(t)\) 作为自身的模长度为一，但作为 \(k\) 向量空间无限维。研究有限维表示时，基域维数可以保证有限长度；研究一般环上的模时，应直接核验所需的链条件。

\ProseParagraphBreak
长度不仅取决于底层加法群，还取决于允许哪些标量作用。在\secref{b2:sec:1201}的三角矩阵环中，上右块的左作用使用 \(K\) 标量，右作用只使用 \(k\) 标量；将 \(K/k\) 改为有限扩张后，这个差别可以直接反映在左右正则模的长度上。

\begin{proposition}\label{b2:prop:triangular-left-right-length}
设 \(k\subseteq K\) 为有限域扩张，\(d=[K:k]\)，并取结合含幺环
\[
R=\left\{\,
\begin{pmatrix}a&u\\0&b\end{pmatrix}
\ \middle|\ a,u\in K,\ b\in k
\,\right\}.
\]
幺性左正则模 \({}_RR\) 与幺性右正则模 \(R_R\) 都有限长度，且
\[
\ell({}_RR)=3,\qquad \ell(R_R)=d+2.
\]
这里右模长度按相应的幺性左 \(R^{\mathrm{op}}\) 模的合成列定义。特别地，\(d>1\) 时，同一个环的左右正则模长度不同。
\end{proposition}
\begin{proof}
记矩阵单位 \(e_{11}=\begin{psmallmatrix}1&0\\0&0\end{psmallmatrix}\)、\(e_{22}=\begin{psmallmatrix}0&0\\0&1\end{psmallmatrix}\)，令
\[
j(u)=\begin{pmatrix}0&u\\0&0\end{pmatrix},\qquad J=j(K).
\]
对 \(r=\begin{psmallmatrix}a&v\\0&b\end{psmallmatrix}\in R\)，有 \(rj(u)=j(au)\)、\(j(u)r=j(ub)\)，所以 \(J\) 是双边理想。

按列分解有 \({}_RR=Re_{11}\oplus Re_{22}\)。\(Re_{11}\) 由左上坐标组成，其左作用为 \(K\) 上的标量乘法；\(J\) 上的左作用同样如此。这两个左模的任意非零元素都能用 \(K\) 标量生成全模，所以都简单。\(Re_{22}/J\) 只保留右下坐标，左作用为 \(k\) 上的标量乘法，故它也简单。因此 \(0\subsetneq J\subsetneq Re_{22}\) 是两层合成列，再由\cref{b2:prop:length-additive}得
\[
\ell({}_RR)=\ell(Re_{11})+\ell(Re_{22})=1+2=3.
\]

按行分解则有 \(R_R=e_{11}R\oplus e_{22}R\)。\(e_{22}R\) 只保留右下坐标，是按 \(k\) 标量右乘的简单右模。又 \(J\subseteq e_{11}R\)，商 \(e_{11}R/J\) 只保留左上坐标，其右作用是乘 \(a\in K\)；任意非零坐标都能生成整个 \(K\)，所以这个商也简单。

\(J\) 上的右作用却只用 \(b\in k\)。选 \(K/k\) 的一组基 \(\omega_1,\ldots,\omega_d\)，令 \(J_0=0\)、\(J_i=j(\operatorname{span}_k\{\omega_1,\ldots,\omega_i\})\)。于是
\[
0=J_0\subsetneq J_1\subsetneq\cdots\subsetneq J_d=J
\]
是右子模链，每个商 \(J_i/J_{i-1}\) 都是一维 \(k\) 空间，且右作用为右下坐标的任意 \(k\) 标量，故都简单。因此 \(\ell(J_R)=d\)。将这条列后接简单商 \(e_{11}R/J\)，得到 \(\ell(e_{11}R)=d+1\)；将右模视为幺性左 \(R^{\mathrm{op}}\) 模，再应用\cref{b2:prop:length-additive}，便得
\[
\ell(R_R)=\ell(e_{11}R)+\ell(e_{22}R)=d+1+1=d+2.
\]
\end{proof}

% !TeX root = ../../Book2.tex
% 纲要标题：### 12.3 不可分解模
% 纲要位置：工作记录/计划纲要.md:L1448

\section{不可分解模}
\label{b2:sec:1203}

合成列逐层取商，直和分解则要求每个分量在原模中有一个补模。两者的差别可由投影看清：一个直和分量对应一个满足 \(e^2=e\) 的自同态。本节先用自同态环描述不可分解性，再证明有限长度条件如何使模分解为自同态作用的幂零分量与可逆分量。

% 纲要要点（供目录核对）：
%
% * 局部自同态环的判别
% * Fitting 引理
% * 有限长度模上的幂等元分解
%

\subsection{局部自同态环的判别}
\label{b2:subsec:1203:01}

\begin{definition}\label{b2:def:indecomposable-module}
设 \(R\) 为含幺环，\(M\) 为非零幺性左 \(R\) 模。若不存在两个非零子模 \(A,B\) 使 \(M=A\oplus B\)，则称 \(M\) 为\term[bukefenjie-mo]{不可分解模}[indecomposable module]。即每次内部直和分解中，至少一个分量必须为零。
\end{definition}

简单模一定不可分解，但逆命题不成立。例如对素数 \(p\)，\(M=\mathbb Z/p^2\mathbb Z\) 的每个非零子群都包含唯一的阶 \(p\) 子群，所以两个非零子模不可能交于零。因此它不可分解，却不是简单模。这里并不否认其有两层合成列，而是说这两层不能在原模中分裂成直和。

\begin{corollary}\label{b2:prop:idempotent-direct-sum}
设 \(R\) 为含幺环，\(M\) 为非零幺性左 \(R\) 模。\(M\) 不可分解，当且仅当 \(\operatorname{End}_R(M)\) 中的幂等元只有 \(0,1\)。
\end{corollary}
\begin{proof}
\cref{b2:prop:idempotent-module-decomposition}已说明，幂等元 \(e\) 给出 \(M=\operatorname{im}e\oplus\ker e\)。若 \(e\ne0\)，像非零；若核为零，则 \(e\bigl(m-e(m)\bigr)=0\) 迫使 \(m=e(m)\) 对所有 \(m\) 成立，即 \(e=1\)。因此 \(e\ne0,1\) 恰好给出两个非零分量。反过来，两个非零分量的坐标投影显然既不为零，也不为恒等。
\end{proof}

比较两个直和分解时，会在某个分量的自同态环中得到恒等映射的有限和表达。若能从这个和中找出一个可逆项，就能比较它与另一分解中的某个分量。为此，我们需要一种环，使有限个非单位相加仍为非单位；这样它们的和便不可能是一。

\begin{definition}\label{b2:def:local-ring}
设 \(E\) 为非零含幺环。若 \(E\) 的全部非单位组成一个双边理想，则称 \(E\) 为\term[jubu-huan]{局部环}[local ring]。这里“单位”指具有两侧逆元的元素；不预设 \(E\) 交换。
\end{definition}

这里的局部性由环的非单位决定，与第六章的局部化构造是不同概念；局部化是一种使选定元素可逆的构造，并非每次局部化所得的环都是局部环。文献也常用“具有唯一极大左理想”定义局部环；在非交换情形，它与上述定义以及唯一极大右理想的条件等价。

\begin{proposition}\label{b2:prop:local-ring-criterion}
设 \(E\) 为非零含幺环。以下条件等价：\(E\) 为局部环；对每个 \(a\in E\)，\(a\) 与 \(1-a\) 中至少有一个为单位；\(E\) 有唯一极大左理想；\(E\) 有唯一极大右理想。在这些条件下，非单位组成的双边理想 \(J\) 是唯一极大左理想，也是唯一极大右理想，\(E/J\) 为除环。有限个非单位之和仍为非单位，特别地，若 \(1=x_1+\cdots+x_r\)，其中 \(r\ge1\)、\(x_i\in E\)，则至少一个 \(x_i\) 为单位。\(E\) 的幂等元只有 \(0,1\)。因此，若 \(R\) 为含幺环、\(M\) 为非零幺性左 \(R\) 模且 \(\operatorname{End}_R(M)\) 为局部环，则 \(M\) 不可分解。
\end{proposition}
\begin{proof}
若非单位构成理想 \(J\)，\(a\) 和 \(1-a\) 不能同时在 \(J\)，否则 \(1\in J\)，与一为单位矛盾。因此所述条件成立。

反过来，假设对每个 \(a\in E\)，\(a\) 或 \(1-a\) 为单位。先证明幂等元 \(e\) 只能为零或一：若 \(e\) 可逆，由 \(e^2=e\) 得 \(e=1\)；若 \(1-e\) 可逆，由 \(e(1-e)=0\) 得 \(e=0\)。这个限制还把一侧逆变成两侧逆。若 \(ab=1\)，则
\[
(ba)^2=b(ab)a=ba.
\]
幂等元 \(ba\) 不能为零，否则 \(b=bab=0\)，与 \(ab=1\) 矛盾。因此 \(ba=1\)，\(a,b\) 互为两侧逆。这一计算对只有左逆或只有右逆的情形都适用。

令 \(J\) 为非单位集合。若 \(a\in J\)、\(r\in E\)，\(ra\) 若可逆，则 \((ra)^{-1}r\) 是 \(a\) 的左逆；\(ar\) 若可逆，则 \(r(ar)^{-1}\) 是 \(a\) 的右逆。上段均会迫使 \(a\) 可逆，矛盾。所以 \(ra,ar\in J\)，并且 \(-a\in J\)。若 \(a,b\in J\) 而 \(u=a+b\) 可逆，则 \(u^{-1}a\) 与 \(1-u^{-1}a=u^{-1}b\) 都非单位，与假设矛盾。因此 \(a+b\in J\)，\(J\) 为不含一的双边理想。

对局部环的这个 \(J\)，任意非零剩余类由 \(a\notin J\) 代表，\(a\) 可逆，故 \(E/J\) 为除环。因此 \(J\) 在左右两侧均极大。任何真左理想或真右理想不能包含单位，故必包含于 \(J\)，所以它们各自的极大理想都只能是 \(J\)。

为从唯一极大左理想推出局部性，要用到每个真左理想都包含于一个极大左理想。对给定真左理想 \(L\)，考虑包含 \(L\) 的全部真左理想，按包含关系排序。它们组成非空偏序集；非空链的并仍是左理想，因为有限个元素可放进同一个链成员中检查减法与左乘，而这个并不含一，否则某个成员已经含一。空链以 \(L\) 为上界。由第一卷附录A定理A.1.4的 Zorn 引理，这个偏序集有极大元，即包含 \(L\) 的极大左理想。

现设 \(J\) 是 \(E\) 的唯一极大左理想。若 \(e^2=e\) 且 \(e\ne0,1\)，则 \(Ee\) 与 \(E(1-e)\) 都是真左理想：例如 \(re=1\) 会给出 \(e=(re)e=r e^2=1\)。由刚证的存在性，两者都包含于 \(J\)，便有 \(1=e+(1-e)\in J\)，矛盾。因此幂等元只有 \(0,1\)，前面关于 \(ba\) 的计算再次说明一侧逆必为两侧逆。

于是非单位 \(a\) 使 \(Ea\) 成为真左理想，必有 \(a\in J\)；反之 \(J\) 不含单位，所以 \(J\) 恰为非单位集合。它已对左乘封闭；若 \(a\in J\) 而 \(ar\) 可逆，则 \(a\) 有右逆，从而为单位，矛盾。因此 \(J\) 也对右乘封闭，是双边理想，\(E\) 为局部环。唯一极大右理想的情形在反环中应用同一论证。有限个非单位之和仍在理想 \(J\) 中，故不可能为单位；若它们之和为一，就至少有一项不在 \(J\) 中。幂等元结论与\cref{b2:prop:idempotent-direct-sum}还给出最后的不可分解性。
\end{proof}

例如，若 \(k\) 为域、\(n\ge1\)，则 \(k[t]/(t^n)\) 中元素可逆，当且仅当其常数项非零。确实，常数项非零的元素可写为 \(a(1-u)\)，其中 \(a\in k^\times\)、\(u\in(t)\) 幂零，逆由有限几何和给出；常数项为零的元素都在真理想 \((t)\)，不可能可逆。故非单位恰为 \((t)\)，这个环为局部环。但 \(\mathbb Z\) 不是局部环：非单位 \(2,-3\) 的和为单位 \(-1\)。虽然后者作为自身的模不可分解，反向局部性还需要下面的有限长度条件。

\subsection{Fitting 引理}
\label{b2:subsec:1203:02}

\begin{lemma}[Fitting]\label{b2:lem:fitting}
设 \(R\) 为含幺环，\(M\) 为有限长度幺性左 \(R\) 模，\(f\in\operatorname{End}_R(M)\)。对所有充分大的正整数 \(n\)，有
\[
M=\ker f^n\oplus\operatorname{im}f^n.
\]
在第一分量上 \(f\) 幂零，在第二分量上 \(f\) 可逆；两个分量都在 \(f\) 下稳定。可取任意 \(n\ge\max\bigl(1,\ell(M)\bigr)\)。
\end{lemma}
\begin{proof}
反复施加 \(f\) 产生两条链：
\[
\begin{aligned}
0&=\ker f^0\subseteq\ker f\subseteq\ker f^2\subseteq\cdots,\\
M&=\operatorname{im}f^0\supseteq\operatorname{im}f\supseteq\operatorname{im}f^2\supseteq\cdots.
\end{aligned}
\]
核记录逐次作用后变成零的元素，像记录仍能由反复作用得到的元素。有限长度同时限制严格升链与严格降链，所以两者都最终稳定。若某步 \(\ker f^r=\ker f^{r+1}\)，则对 \(f^{r+2}(x)=0\)，有 \(f(x)\in\ker f^{r+1}=\ker f^r\)，继而 \(x\in\ker f^{r+1}=\ker f^r\)。逐次归纳说明核链从此稳定。令 \(d=\ell(M)\)，从 \(\ker f^0=0\) 到 \(\ker f^d\) 若有一步相等，核链便从该步稳定；若前 \(d\) 步全都严格增加，则 \(\ell(\ker f^d)\ge d\)，所以 \(\ker f^d=M\)，此后也不再增加。因此从 \(r=d\) 起核链必稳定。由\cref{b2:prop:finite-length-endomorphism}，
\[
\ell(\operatorname{im}f^r)=\ell(M)-\ell(\ker f^r),
\]
故像链从同一位置起长度不变。相邻像有包含关系，长度相等时商长度为零，所以像也相等。

取满足上述界的 \(n\)，于是 \(\ker f^{2n}=\ker f^n\)、\(\operatorname{im}f^{2n}=\operatorname{im}f^n\)。要证明核与像交于零，考虑既经过 \(n\) 次作用、又会在 \(n\) 次作用后变成零的元素 \(y\)。写 \(y=f^n(x)\) 且 \(f^n(y)=0\)，就有 \(f^{2n}(x)=0\)。这正是比较 \(2n\) 与 \(n\) 的原因；由核稳定，\(x\in\ker f^{2n}=\ker f^n\)，从而 \(y=0\)。

对任意 \(x\in M\)，还要从像中减去一项，使余项落入核。由像稳定，\(f^n(x)\in\operatorname{im}f^{2n}\)，所以可取 \(z\in M\)，使 \(f^n(x)=f^{2n}(z)\)。减去 \(f^n(z)\) 后有
\[
f^n\bigl(x-f^n(z)\bigr)=0.
\]
因此 \(x-f^n(z)\in\ker f^n\)，而 \(f^n(z)\in\operatorname{im}f^n\)，得到所需分解。若 \(x\in\ker f^n\)，则 \(f^n(f(x))=f(f^n(x))=0\)，所以核在 \(f\) 下稳定；像也稳定，因为 \(f(\operatorname{im}f^n)=\operatorname{im}f^{n+1}=\operatorname{im}f^n\)。在核上 \(f^n=0\)；在像上，这个等式说明 \(f\) 满射，而 \(\ker f\subseteq\ker f^n\) 与已证的零交说明限制映射也单射，故可逆。
\end{proof}

这称为\term[Fitting-yinli]{Fitting 引理}[Fitting lemma]。\(\max(1,\ell(M))\) 是对所有自同态通用的指数上界，不一定是某个 \(f\) 的最小分解指数；\cref{b2:prop:fitting-integer-projections}中模的长度为五，乘二却在第三次作用时就得到分解。

\ProseParagraphBreak
对域 \(k\) 上有限维空间的算子 \(T\)，稳定核由那些被某个 \(T\) 的幂消去的向量组成，正是广义零特征子空间；稳定像是 \(T\) 可逆的部分。若最小多项式在 \(k\) 上分裂，\cref{b2:thm:jordan-canonical-form}的 Jordan 形把所有 \(J_r(0)\) 块放在第一部分，把非零特征值的块放在第二部分。Fitting 引理没有要求底环是域，也没有要求多项式分裂；有限长度提供的是核和像的同时稳定。

\begin{theorem}\label{b2:thm:indecomposable-local-end}
设 \(R\) 为含幺环，\(M\) 为非零幺性左 \(R\) 模。若 \(\operatorname{End}_R(M)\) 为局部环，则 \(M\) 不可分解。若 \(M\) 还有有限长度，则 \(M\) 不可分解，当且仅当 \(\operatorname{End}_R(M)\) 为局部环；在有限长度且不可分解情形，每个 \(f\in\operatorname{End}_R(M)\) 或者可逆，或者幂零，其全部非可逆自同态构成双边理想。
\end{theorem}
\begin{proof}
局部自同态环的幂等元只有零和一，由\cref{b2:prop:idempotent-direct-sum}，\(M\) 不可分解。这一方向没有使用长度假设。现设 \(M\) 有限长度且不可分解。对任意 \(f\)，Fitting 引理给出核与像的直和分解，不可分解性迫使
\[
\ker f^n=0\qquad\text{或}\qquad\operatorname{im}f^n=0.
\]
前一种情形使 \(f\) 单射，由\cref{b2:prop:finite-length-endomorphism}可逆；后一种情形使 \(f^n=0\)，即 \(f\) 幂零。

当 \(f\) 幂零时，\(1-f\) 可逆，因为若 \(f^n=0\)，有
\[
(1-f)(1+f+\cdots+f^{n-1})
=(1+f+\cdots+f^{n-1})(1-f)=1.
\]
因此每个 \(f\) 都使 \(f\) 或 \(1-f\) 可逆。\cref{b2:prop:local-ring-criterion}给出自同态环为局部环，同时保证非可逆元素对加法及左右乘法封闭。
\end{proof}

局部环一般只保证非单位组成理想，并不保证每个非单位都幂零。这里的幂零结论同时使用了有限长度与不可分解性；\cref{b2:prop:local-ring-nonnilpotent}将给出局部环中非幂零非单位的例子。

\ProseParagraphBreak

单有有限长度也不足以保证非可逆自同态幂零。例如对域 \(k\)，\(k^2\) 上的投影 \(\operatorname{diag}(1,0)\) 不可逆，也不幂零。相反，对循环 \(k[t]\) 模 \(k[t]/(p^n)\)，其中 \(p\) 不可约、\(n\ge1\)，任意两个非零子模均含有最小非零子模 \((p^{n-1})/(p^n)\)，所以它不可分解；其非可逆自同态因而都幂零。这与其自同态环 \(k[t]/(p^n)\) 中的非单位为 \((p)/(p^n)\) 相吻合。

\subsection{有限长度模上的幂等元分解}
\label{b2:subsec:1203:03}

\cref{b2:prop:idempotent-module-decomposition}也给出了有限版本：若自同态满足
\[
e_i^2=e_i,\qquad e_ie_j=0\ (i\ne j),\qquad
\sum_{i=1}^re_i=1,
\]
则 \(M=\bigoplus_i e_iM\)，反过来有限内部直和的坐标投影满足这些等式。现在进一步判别一个投影的像是否还能分解。

\begin{proposition}\label{b2:prop:orthogonal-idempotent-decomposition}
设 \(R\) 为含幺环，\(M\) 为幺性左 \(R\) 模。对非零幂等元 \(e\in\operatorname{End}_R(M)\)，其像 \(eM\) 不可分解，当且仅当 \(e\) 不能写成两个非零幂等元 \(u,v\) 之和且 \(uv=vu=0\)。若 \(M\) 有限长度，\(I\) 为集合，且 \((e_i)_{i\in I}\) 是 \(\operatorname{End}_R(M)\) 中满足 \(e_i\ne0\)、\(e_i^2=e_i\) 及 \(e_ie_j=0\ (i\ne j)\) 的一族元素，则 \(I\) 有限且 \(|I|\le\ell(M)\)。
\end{proposition}
\begin{proof}
若 \(e=u+v\) 如题述，则 \(eu=ue=u\)、\(ev=ve=v\)，所以两像都在 \(eM\) 中，并有 \(eM=uM\oplus vM\)，两项非零。反过来，若 \(eM=A\oplus B\) 且两项非零，先用 \(e\) 把 \(M\) 投到 \(eM\)，再投到 \(A,B\)，得到 \(u,v\)。它们分别在自身像上为恒等，在另一像及 \(\ker e\) 上为零，因而幂等、正交、非零，且 \(u+v=e\)。

对正交幂等元族，只需先看任意有限子集 \(F\subseteq I\)。若 \(x_i\in e_iM\) 且 \(\sum_{i\in F}x_i=0\)，施加 \(e_j\) 得 \(x_j=0\)，因为 \(e_j\) 在 \(e_jM\) 上为恒等，而在其余各项上为零。因此 \(\sum_{i\in F}e_iM\) 是内部直和。每个 \(e_i\ne0\) 使 \(e_iM\ne0\)，长度至少一，故由长度可加性，
\[
|F|\le\sum_{i\in F}\ell(e_iM)
=\ell\!\left(\bigoplus_{i\in F}e_iM\right)
\le\ell(M).
\]
若 \(I\) 有多于 \(\ell(M)\) 个元素，就能取到一个含 \(\ell(M)+1\) 个元素的有限子集，与此矛盾。所以 \(I\) 有限且满足所述界。整个论证只使用有限子族，没有对无限族定义幂等元之和。
\end{proof}

“正交”在这里指乘积为零，不涉及内积。若 \(M\) 有限长度，每个非零投影的像至少占一层长度，因此不能无限细分投影。如果 \(eM=A\oplus B\) 且两项非零，则 \(\ell(eM)=\ell(A)+\ell(B)\)，两个新分量的长度都严格小于原分量；继续分解时可对长度归纳。第十五章将把不能再写成两个非零正交幂等元之和的幂等元称为本原幂等元。

\ProseParagraphBreak
对内部直和 \(M=A\oplus B\)，自同态可写成由四个 Hom 群组成的分块矩阵；投影到 \(A\) 的幂等元是 \(\begin{pmatrix}1&0\\0&0\end{pmatrix}\)。在一般情形，它不在整个自同态环的中心，因为与含有非零交叉映射的分块矩阵不一定交换。模的直和分解需要的是幂等元，并不要求这些投影是中心幂等元；中心性对应更强的环论分解，后面另行讨论。

\begin{proposition}\label{b2:prop:fitting-integer-projections}
把 \(M=\mathbb Z/72\mathbb Z\) 看作幺性左 \(\mathbb Z\) 模，令 \(f(\bar x)=2\bar x\)。使 \(M=\ker f^n\oplus\operatorname{im}f^n\) 成立的最小正整数为 \(n=3\)，此时
\[
\ker f^3=\langle\bar9\rangle,\qquad
\operatorname{im}f^3=\langle\bar8\rangle.
\]
两个分量的投影分别为乘九与乘六十四，\(f\) 在第二分量上的逆为乘五。\(M\) 的长度为五。
\end{proposition}
\begin{proof}
\(f^3\) 为乘八，而 \(72=8\cdot9\)，故 \(8x\equiv0\pmod{72}\) 等价于 \(9\mid x\)。这给出所列核；像由 \(\bar8\) 生成。两分量分别有八与九个元素，交集为零，因为同时被八和九整除的整数被七十二整除。又 \(\bar9-\bar8=\bar1\)，所以它们共同生成 \(M\)，得到直和分解。这也就是中国剩余分解 \(\mathbb Z/72\mathbb Z\cong\mathbb Z/8\mathbb Z\oplus\mathbb Z/9\mathbb Z\)：乘二在前者上幂零，在后者上可逆。

乘九在 \(\langle\bar9\rangle\) 上为恒等，因为 \(9^2\equiv9\pmod{72}\)，在 \(\langle\bar8\rangle\) 上为零，因为 \(9\cdot8\equiv0\pmod{72}\)。因此它是第一分量的投影；另一分量的投影为 \(1-9=-8\equiv64\pmod{72}\)，即乘六十四。在阶九的第二分量上，乘二与乘五的复合为乘十，而 \(10\equiv1\pmod9\)，故乘五为所需逆。

当 \(n=1\) 时，\(\ker f=\langle\overline{36}\rangle\)，其中非零的 \(\overline{36}\) 也在 \(\operatorname{im}f=2M\) 中；当 \(n=2\) 时，\(\ker f^2=\langle\overline{18}\rangle\)，\(\overline{36}\) 又在 \(\operatorname{im}f^2=4M\) 中。因此前两步的核与像交集非零，最小指数确为三。最后，
\[
0\subsetneq\langle\overline{36}\rangle
\subsetneq\langle\overline{18}\rangle
\subsetneq\langle\bar9\rangle
\subsetneq\langle\bar3\rangle
\subsetneq M
\]
的相邻商依次为阶二、二、二、三、三的循环群，都是简单 \(\mathbb Z\) 模，所以这是长度五的合成列。
\end{proof}

\begin{proposition}\label{b2:prop:local-ring-nonnilpotent}
固定素数 \(p\)，令
\[
R=\mathbb Z_{(p)}
=\bigl\{\,a/b\in\mathbb Q\bigm|a,b\in\mathbb Z,\ p\nmid b\,\bigr\}.
\]
则 \(R\) 为交换局部环，非单位理想为 \(pR\)。其正则幺性左模 \({}_RR\) 不可分解，却不满足降链条件，因而不具有有限长度。非单位 \(p\in R\) 以及正则模上的乘 \(p\) 自同态都不幂零。
\end{proposition}
\begin{proof}
分母不被 \(p\) 整除的分数对加法、减法与乘法封闭，并含零与一，所以 \(R\) 是 \(\mathbb Q\) 的含幺子环，特别地为整环。对 \(a/b\in R\)，若 \(p\nmid a\)，其逆 \(b/a\) 仍在 \(R\) 中；若 \(p\mid a\)，则它在 \(R\) 中不可能有逆。否则存在 \(c/d\in R\) 满足 \(ac=bd\)，但左边被 \(p\) 整除、右边不被 \(p\) 整除，矛盾。因此非单位恰为 \(pR\)，这个集合是双边真理想，\(R\) 为局部环。由\cref{b2:prop:regular-hom-endomorphism}，\(\operatorname{End}_R({}_RR)\cong R^{\mathrm{op}}=R\)，故\cref{b2:prop:local-ring-criterion}推出正则模不可分解。

正则模的子模降链
\[
R\supsetneq pR\supsetneq p^2R\supsetneq\cdots
\]
处处严格：若 \(p^n\in p^{n+1}R\)，在分式域中消去 \(p^{n+1}\) 就得到 \(1/p\in R\)，矛盾。所以正则模不是 Artin 模，也不具有有限长度。由于 \(R\subseteq\mathbb Q\)，\(p^n\ne0\) 对每个 \(n\ge1\) 都成立；乘 \(p\) 自同态的第 \(n\) 次幂在元素一上取值为 \(p^n\)，同样非零。这给出局部自同态环中不幂零的非单位，而有限长度假设在这里恰好失败。
\end{proof}

% !TeX root = ../../Book2.tex
% 纲要标题：### 12.4 Krull–Schmidt 定理
% 纲要位置：工作记录/计划纲要.md:L1454

\section{Krull–Schmidt 定理}
\label{b2:sec:1204}

有限长度使一次非平凡直和分解的每个分量都变短，因而分解过程会终止。唯一性则更细：两个分解中的分量未必是同一个子模，要先证明某一对分量同构，再交换相应的补模，才能继续比较。局部自同态环恰好保证这一交换能够进行。

% 纲要要点（供目录核对）：
%
% * 有限长度情形的不可分解分解存在性
% * 局部自同态环、补模交换与不可分解因子的唯一性
% * 无限乘积的存在性失效及 Noether 模的非唯一分解
%
% ---
%

\subsection{有限长度模的分解存在性}
\label{b2:subsec:1204:01}

\begin{proposition}\label{b2:prop:finite-length-decomposition}
设 \(R\) 为含幺环，\(M\) 为有限长度幺性左 \(R\) 模。若 \(M\ne0\)，则 \(M\) 是有限多个非零不可分解左 \(R\) 模的直和，分量个数不超过 \(\ell(M)\)。零模对应空直和。
\end{proposition}
\begin{proof}
零模取空直和。对非零模的长度 \(\ell(M)\) 归纳。长度一时 \(M\) 简单，因而不可分解。一般地，若 \(M\) 不可分解，取它本身为唯一分量。否则可写成 \(M=A\oplus B\)，其中两项非零。由\cref{b2:prop:length-additive}，\(\ell(M)=\ell(A)+\ell(B)\)，所以两项长度都是严格小于 \(\ell(M)\) 的正整数，可以分别使用归纳假设，再将两个有限分解合并。沿每次继续分解的分量，长度都严格下降，因而分解过程不能无限继续。每个非零分量长度至少一，长度可加性还给出分量个数的上界。
\end{proof}

例如，若 \(k\) 为域，则 \(t^2\) 与 \((t-1)^3\) 互素，\(k[t]\) 模 \(k[t]/\bigl(t^2(t-1)^3\bigr)\) 通过\secref{b2:sec:0403}证明的中国剩余同构分成
\[
k[t]/(t^2)\oplus k[t]/\bigl((t-1)^3\bigr).
\]
这正是\cref{b2:thm:pid-elementary-divisors}中的两个循环主分量，分别对应不可约多项式 \(t\) 与 \(t-1\)。每个循环商的非零子模都含有其最小非零子模，分别为 \((t)/(t^2)\) 与 \(\bigl((t-1)^2\bigr)/\bigl((t-1)^3\bigr)\)，所以两个分量均不可分解，长度分别为二和三。它们各自仍有非平凡合成列：不可分解性不等于没有真子模。

\ProseParagraphBreak
分解存在性也可以表述为投影的逐步细分。若一个分量还能分解，就把对应幂等元拆成两个非零正交幂等元。\cref{b2:prop:orthogonal-idempotent-decomposition}说明这与模的操作一致；长度每次被分到两个正整数中，所以不可能无限继续。这个证明保证存在有限分解，却尚未比较不同选择得到的分量。

\subsection{不可分解因子的唯一性}
\label{b2:subsec:1204:02}

\begin{theorem}[Krull–Schmidt]\label{b2:thm:krull-schmidt}
设 \(R\) 为含幺环，\(M\) 为有限长度幺性左 \(R\) 模。若 \(M\) 有两个分解
\[
M=A_1\oplus\cdots\oplus A_r
=B_1\oplus\cdots\oplus B_s,
\]
其中各项非零且不可分解，则 \(r=s\)，并可重排下标，使 \(A_i\cong B_i\) 对所有 \(i\) 成立。因此不可分解直和分量的同构类型及重数由 \(M\) 唯一确定。
\end{theorem}
\begin{proof}
对 \(\ell(M)\) 归纳，零模只有空分解。非零时，令
\[
a_i:A_i\rightarrow M,\quad p_i:M\rightarrow A_i,\qquad
b_j:B_j\rightarrow M,\quad q_j:M\rightarrow B_j
\]
为两组包含和投影。复合 \(p_1b_jq_ja_1\) 先把 \(A_1\) 中的元素包含到 \(M\)，取出它的第 \(j\) 个 \(B\) 分量，将这个分量送回 \(M\)，最后取出其中的 \(A_1\) 分量。所有这些往返映射相加，应当恢复原元素：由 \(\sum_jb_jq_j=\operatorname{id}_M\) 和 \(p_1a_1=\operatorname{id}_{A_1}\)，在 \(\operatorname{End}_R(A_1)\) 中有
\begin{equation}\label{b2:eq:krull-schmidt-unit-sum}
\operatorname{id}_{A_1}=\sum_{j=1}^s p_1b_jq_ja_1.
\end{equation}
\(A_1\) 是非零不可分解有限长度模，故这个自同态环由\cref{b2:thm:indecomposable-local-end}为局部环。右侧只有有限项，左侧恒等自同态是单位；\cref{b2:prop:local-ring-criterion}中的有限和性质保证至少一项可逆。重排后设 \(p_1b_1q_1a_1\) 可逆。

记 \(u=q_1a_1:A_1\rightarrow B_1\)、\(v=p_1b_1:B_1\rightarrow A_1\)，于是 \(vu\) 可逆。要把 \(u(A_1)\) 分裂成 \(B_1\) 的直和因子，需要给 \(u\) 构造一个左逆。原映射 \(v\) 与 \(u\) 的复合为可逆自同态 \(vu\)，因此先用 \((vu)^{-1}\) 修正它，令
\[
w=(vu)^{-1}v:B_1\longrightarrow A_1.
\]
便有 \(wu=\operatorname{id}_{A_1}\)，所以 \(u\) 单射。幂等元 \(uw\) 在 \(B_1\) 上的像为 \(u(A_1)\)、核为 \(\ker w\)，由\cref{b2:prop:idempotent-module-decomposition}，
\[
B_1=u(A_1)\oplus\ker w.
\]
由于 \(A_1\ne0\)，第一项非零；\(B_1\) 不可分解，故 \(\ker w=0\)，上式只剩 \(B_1=u(A_1)\)。因此 \(u\) 满且单，得到 \(A_1\cong B_1\)。这里起作用的是：非零模分裂嵌入不可分解模后，其像必须占据整个目标模。

虽然 \(A_1\cong B_1\)，它们在 \(M\) 中的位置仍可能不同，尚不能据此约去两边的同构项。令 \(B'=\bigoplus_{j>1}B_j\)，要把实际子模 \(A_1\) 换入第二个分解，同时保留同一个补模 \(B'\)。因为 \(q_1|_{A_1}=u\) 为同构，对任意 \(m\in M\)，它的 \(B_1\) 坐标 \(q_1(m)\) 都对应唯一的 \(a\in A_1\)。于是 \(q_1(a)=q_1(m)\)，从而 \(m-a\in\ker q_1=B'\)，所以 \(A_1+B'=M\)。若 \(a\in A_1\cap B'\)，则 \(u(a)=0\)，所以 \(a=0\)。这就得到实际的直和分解 \(M=A_1\oplus B'\)。

令 \(A'=\bigoplus_{i>1}A_i\)。两个实际分解 \(M=A_1\oplus A'=A_1\oplus B'\) 都将其补模同构到商 \(M/A_1\)：限制商映射的核为零，且每个陪集都有补模中的代表。因此
\[
A'\cong M/A_1\cong B'.
\]
将 \(B'\) 的分解沿此同构转移到 \(A'\)，就得到同一个模上的两个不可分解分解。由长度可加性，\(\ell(A')=\ell(M)-\ell(A_1)<\ell(M)\)，故可用归纳假设，得到剩余项一一同构。加上 \(A_1\cong B_1\)，完成证明。
\end{proof}

这一\term[Krull-Schmidt-dingli]{Krull–Schmidt 定理}[Krull–Schmidt theorem]的唯一性在同构和排列意义下成立。分量作为子模的具体位置可以不同，例如 \(k^2=ke_1\oplus ke_2=ke_1\oplus k(e_1+e_2)\)，第二个分量改变了；两种分解的不可分解因子都为两个一维空间。

\begin{corollary}[有限长度模的消去]\label{b2:cor:finite-length-cancellation}
设 \(R\) 为含幺环，\(A,B,C\) 为有限长度幺性左 \(R\) 模。若 \(A\oplus C\cong B\oplus C\)，则 \(A\cong B\)。
\end{corollary}
\begin{proof}
将 \(A,B,C\) 分别分成有限个不可分解模。由\cref{b2:thm:krull-schmidt}，左右两侧每一种不可分解同构类型的出现次数相等。两侧减去来自 \(C\) 的相同有限重数，得到 \(A,B\) 的分量类型及重数相同，逐项同构的直和便给出 \(A\cong B\)。
\end{proof}

Jordan–Hölder 定理比较的是子模链中的简单商层，Krull–Schmidt 定理比较的是原模内部的不可分解直和因子。回看\subsecref{b2:subsec:1202:01}中比较过的 \(\mathbb Z/p^2\mathbb Z\) 与 \((\mathbb Z/p\mathbb Z)^2\)，其中 \(p\) 为素数：二者的合成因子都是两个 \(\mathbb Z/p\mathbb Z\)，但前者的任意非零子模都包含 \(p\mathbb Z/p^2\mathbb Z\)，所以自身不可分解；后者则已经分成两个简单直和因子。相同的简单商层不能说明它们怎样拼合，而不可分解分解会区分这两种拼合。对有限维代数的表示，二者的有限维版本可参阅 Etingof 等人\cite[Theorem 3.7.1, pp. 50--51; Theorem 3.8.1, pp. 51--52; Remark 3.8.6, p. 53]{EtingofEtAl2011RepresentationTheory}；本章用链条件与 Fitting 引理处理一般有限长度模。

\subsection{分解定理的条件与反例}
\label{b2:subsec:1204:03}

有限长度在证明中有两个作用。分解存在性用到的是非平凡分量的正整数长度严格下降，迫使分解过程终止；唯一性则用到核、像的两种链同时稳定，由 Fitting 引理得到不可分解模的自同态环局部，再从恒等映射的有限和中找到可逆项。去掉有限长度以后，这两步都可能失败，甚至有限生成或 Noether 条件也不足以替代它。

\ProseParagraphBreak
先看自同态的失效。整数模 \(\mathbb Z\) 不可分解，因为两个非零子群总有非零交；其自同态环却是非局部环 \(\mathbb Z\)。乘二的映射单射而不满，任意 \(n\ge1\) 都有 \(\ker(2^n)=0\)、\(\operatorname{im}(2^n)=2^n\mathbb Z\ne\mathbb Z\)，所以 Fitting 分解也不成立。缺少降链条件使像永远不能稳定。这个模已有以自身为唯一分量的不可分解分解；此处失去的是从不可分解性推出自同态环局部的机制。

\begin{example}[没有不可分解直和分解的循环模]\label{b2:ex:no-indecomposable-decomposition}
设 \(k\) 为域，令含幺环 \(R=\prod_{n\ge1}k\)，按坐标运算，单位元为 \((1,1,\ldots)\)。幺性左正则模 \(R\) 虽由一生成，却不能写成任何一族非零不可分解子模的内部直和。
\end{example}
\begin{proof}
任意直和因子的投影 \(p:R\rightarrow R\) 由 \(e=p(1)\) 决定，\(p(r)=re\)。投影条件给出 \(e^2=e\)，每个坐标因 \(k\) 为域只能是零或一。因此直和因子为 \(eR\)，对应一个坐标子集 \(S\)。

若 \(S\) 至少有两个元素，可分成两个非空不交子集，相应指示幂等元把 \(eR\) 分成两个非零直和因子。若 \(S\) 只有一个元素，\(eR\cong k\) 没有非零真子模，故不可分解。因此每个非零不可分解直和因子只能占据一个坐标。

若 \(R\) 是这些因子的内部直和，即使因子族无限，\cref{b2:prop:internal-direct-sum}也要求每个元素只有有限多个非零分量。因此单位元 \((1,1,\ldots)\) 必须是有限多个分量元素之和。这种和只能有有限个非零坐标，不可能等于单位元，矛盾。
\end{proof}

唯一性也可能真正失败，而不只是上述证明失去效力。以下例子给出两个不同的有限不可分解分解。

\begin{example}[Noether 模的非唯一分解]\label{b2:ex:krull-schmidt-nonunique}
令 \(s=\sqrt{-5}\)、\(R=\mathbb Z[s]\subseteq\mathbb C\)，\(I=(2,1+s)\)。以通常乘法为作用，将 \(R\) 与 \(I\) 视为幺性左 \(R\) 模。它们是不同构的不可分解 Noether 模，但
\[
R\oplus R\cong I\oplus I.
\]
\end{example}
\begin{proof}
先说明 \(I\) 是非零真理想。将整数模二、将 \(s\) 送到一，给出满同态 \(R\rightarrow\mathbb F_2\)，因为关系 \(s^2+5=0\) 在目标中成立。它将 \(2,1+s\) 送到零，故 \(I\) 真；又 \(2\in I\)，所以 \(I\ne0\)。

对 \(a+bs\in R\)，令 \(N(a+bs)=a^2+5b^2\)，这是乘共轭所得的非负整数，满足 \(N(uv)=N(u)N(v)\)。若 \(I=(u)\)，则 \(u\) 同时整除 \(2\) 和 \(1+s\)，故正整数 \(N(u)\) 同时整除四和六，只能为一或二。若为一，共轭元素就是 \(u\) 的逆，\(I=R\)，矛盾；若为二，方程 \(a^2+5b^2=2\) 无整数解。因此 \(I\) 非主，不能与 \(R\) 同构：一个同构会把一送到 \(I\) 的单个生成元。

为构造 \(R^2\) 与 \(I\oplus I\) 之间的同构，需要控制 \(I\) 中元素的乘积。这里 \(I^2\) 表示理想乘积，即由所有 \(xy\)、\(x,y\in I\) 生成的理想；\(I\oplus I\) 则以有序对为元素。计算理想 \(I^2\) 时，生成元 \(4,2(1+s),(1+s)^2=-4+2s\) 都在 \(2R\) 中。另一方面，\(s-1=(1+s)-2\in I\)，并且
\[
2=2(2\cdot2)+(1+s)(s-1),
\]
所以 \(2\in I^2\)，得到 \(I^2=2R\)。现在选一个条目都在 \(I\) 中的 \(2\times2\) 矩阵 \(B\)，它就把 \(R^2\) 送入 \(I\oplus I\)。二维伴随矩阵的条目仍在 \(I\) 中，故 \(\operatorname{adj}(B)\) 乘上 \(I\oplus I\) 中的向量后，分子各坐标都是两个 \(I\) 中元素的乘积之和，属于理想 \(I^2=2R\)。为了让逆矩阵 \((\det B)^{-1}\operatorname{adj}(B)\) 把这些向量送回 \(R^2\)，可以选择 \(\det B=2\)：这样分子中的二恰好抵消分母，所得坐标仍在 \(R\) 中。

为找到这样的 \(B\)，把刚才的 \(2=8+(s-1)(s+1)\) 组织成一个行列式：
\[
B=\begin{pmatrix}2&-(s-1)\\1+s&4\end{pmatrix},\qquad
\det B=2\cdot4+(s-1)(s+1)=2.
\]
\cref{b2:prop:adjugate-identity}给出
\[
B^{-1}=\frac12\operatorname{adj}(B)
=\frac12\begin{pmatrix}4&s-1\\-(1+s)&2\end{pmatrix}.
\]
于是对 \(x,y\in I\)，得到以下公式：
\[
\Phi(x,y)=
\left(2x+\frac{s-1}{2}y,\ -\frac{1+s}{2}x+y\right).
\]
这里的分式系数位于分式域中；但 \((s-1)y\) 与 \((1+s)x\) 都在 \(I^2=2R\) 中，所以两个除以二的表达确实属于 \(R\)，\(\Phi\) 取值于 \(R^2\)。
反向公式为
\[
\Psi(r,z)=\bigl(2r-(s-1)z,\ (1+s)r+4z\bigr).
\]
它确实取值于 \(I\oplus I\)，因为四个系数 \(2,s-1,1+s,4\) 均在 \(I\)。在分式域中写成矩阵，
\[
\begin{pmatrix}2&(s-1)/2\\-(1+s)/2&1\end{pmatrix}
\begin{pmatrix}2&-(s-1)\\1+s&4\end{pmatrix}
=\begin{pmatrix}1&0\\0&1\end{pmatrix};
\]
反序相乘也为单位矩阵，使用的恒等式为 \((s-1)(s+1)=-6\)。两映射都是分式域线性映射在相应子模上的限制，故为 \(R\) 线性；正反公式又已核对实际值域，因而给出 \(I\oplus I\cong R^2\) 的 \(R\) 模同构。

任意两个非零 \(R\) 子模 \(A,B\) 若都在 \(R\) 或 \(I\) 内，取非零 \(a\in A,b\in B\)。由于 \(a,b\) 同时属于标量环 \(R\)，有 \(ab=b\cdot a\in A\)、\(ab=a\cdot b\in B\)，且 \(R\) 为整环使 \(ab\ne0\)。因此这两个子模总有非零交，\(R,I\) 均不可分解。最后，\(R\) 作为整数模自由，基为 \(1,s\)。任意 \(R\) 子模都是 \(\mathbb Z^2\) 的子模，由\cref{b2:thm:pid-free-submodule}有限生成；一组整数生成元同时也是 \(R\) 生成元。所以 \(R\) 为 Noether 模，\(I\) 作为其子模也为 Noether 模。得到的两个分解均有限，却没有逐项同构的可能。
\end{proof}

这里的 \(R\) 与 \(I\) 都不满足 Artin 条件：\(R\supsetneq2R\supsetneq4R\supsetneq\cdots\) 严格下降；又因 \(2\in I\) 而 \(2\notin2I\)（否则 \(1\in I\)），\(I\supsetneq2I\supsetneq4I\supsetneq\cdots\) 也严格下降。因此这个反例不与有限长度情形的定理矛盾。

\ProseParagraphBreak

无限乘积环的循环正则模说明有限生成不保证不可分解分解的存在；\(R,I\) 的例子说明，即使 Noether 模已经有有限不可分解分解，其因子的同构类型也可能不唯一。Krull–Schmidt 还有适用于某些更大范围模的版本，但必须另外给出保证直和分量局部性和分解存在性的假设。


\begin{chaptersummary}
Noether 条件等价于每个子模有限生成，Artin 条件则限制降链；两者都能通过短正合列的两端判定。对非交换环，左理想与右理想给出的条件可能不同。两种链条件同时成立，才保证存在有限合成列；Jordan–Hölder 定理使其层数与简单因子重数成为模的不变量，长度可加性则把它们用于子模、商模和同态计算。

\ProseParagraphBreak
合成列并不自动分裂。不可分解只排除非平凡直和，不排除真子模或非分裂扩张；Fitting 引理使有限长度不可分解模的每个自同态或者可逆、或者幂零，从而使自同态环局部。Krull–Schmidt 唯一性利用这个局部性，从一个恒等自同态的有限和中找出可逆项，再交换对应的直和分量。结论确定不可分解类型和重数，不确定分量的具体嵌入位置。

\ProseParagraphBreak
无限乘积环的正则模说明，循环模也未必存在不可分解直和分解；\(\mathbb Z[\sqrt{-5}]\) 上的理想例子则说明，Noether 模可以有两个不同的有限不可分解分解。这些限制促使我们继续考察哪些环能使所有模或某类模的分解更简单。\chapref{b2:ch:13}的半单模把每个分量进一步要求为简单模，随后再由环的结构解释这种现象。
\end{chaptersummary}

\BookExercises

\begin{exercise}\label{b2:ex:12-vector-chain-conditions}
设 \(k\) 为域，\(V\) 为 \(k\) 向量空间。证明 \(V\) 作为 \(k\) 模满足升链条件、满足降链条件、有限维三者等价。说明仅将“向量空间”替换成“一般环上的模”后，哪些蕴含不再成立。
\end{exercise}
\begin{solution}
有限维时，严格包含使维数至少改变一，所以升链和降链都稳定。若维数无限，沿用本书的选择公理约定，逐次选择不在此前向量生成空间内的向量，得到线性无关序列 \(e_1,e_2,\ldots\)，其前 \(n\) 项的线性生成空间严格上升，其尾部 \(\operatorname{span}_k\{e_i:i\ge n\}\) 严格下降，所以两种链条件都失败。一般模上，\(\mathbb Z\) 是 Noether 而非 Artin 模，\(C_{p^\infty}\) 是 Artin 而非 Noether 模，故两种链条件不能互推。“有限维”也没有对任意环都适用的同一含义，合适的共同条件是\cref{b2:thm:finite-length-chain-conditions}中的有限长度。
\end{solution}

\begin{exercise}\label{b2:ex:12-directed-cover}
设 \(R\) 为含幺环，\(M\) 为幺性左 \(R\) 模。非空子模族 \((N_i)_{i\in I}\) 中若任意两个成员都包含于某个共同的成员，就称该族向上有向。证明 \(M\) 有限生成，当且仅当每个满足 \(M=\bigcup_iN_i\) 的向上有向子模族都含有等于 \(M\) 的一个成员。这个条件为何不足以直接说明 \(M\) 是 Noether 模？
\end{exercise}
\begin{solution}
若 \(M\) 由有限个 \(x_1,\ldots,x_r\) 生成，各元素分别落在族中某项；反复使用有向性，可找到同时包含它们的 \(N_j\)，故 \(N_j=M\)。反过来，取 \(M\) 的全部有限生成子模，它们的并为 \(M\)，任意两个的和仍有限生成，所以是向上有向族。假设迫使某个有限生成子模就是 \(M\)。这个条件只针对 \(M\) 自身的生成性；Noether 条件还要求每个子模有限生成，二者的区别由\cref{b2:prop:noetherian-submodules-fg}明确给出。
\end{solution}

\begin{exercise}\label{b2:ex:12-generator-extension}
设 \(R\) 为含幺环，\(0\rightarrow A\xrightarrow{i}B\xrightarrow{q}C\rightarrow0\) 为幺性左 \(R\) 模的短正合列。若 \(A\) 可由 \(r\) 个元素生成，\(C\) 可由 \(s\) 个元素生成，证明 \(B\) 可由至多 \(r+s\) 个元素生成，并说明在哪里使用了正合性。
\end{exercise}
\begin{solution}
取 \(A\) 的生成元 \(a_1,\ldots,a_r\)，\(C\) 的生成元 \(c_1,\ldots,c_s\)。由 \(q\) 满射，取 \(b_j\in B\) 使 \(q(b_j)=c_j\)。给定 \(b\in B\)，写 \(q(b)=\sum_j t_jc_j\)，则 \(b-\sum_jt_jb_j\in\ker q=\operatorname{im}i\)，所以又可写成 \(\sum_\nu u_\nu i(a_\nu)\)。于是 \(i(a_\nu),b_j\) 共同生成 \(B\)。满射性用于选取商生成元的原像，\(\ker q=\operatorname{im}i\) 用于把差送回 \(A\)；只知道两个映射可复合，并不足以得到这个结论。
\end{solution}

\begin{exercise}\label{b2:ex:12-one-sided-endomorphism}
设 \(R\) 为含幺环，\(M\) 为幺性左 \(R\) 模。证明 \(M\) 为 Noether 模时每个满射自同态可逆，\(M\) 为 Artin 模时每个单射自同态可逆。分别给出不能互换“满射”与“单射”的例子。
\end{exercise}
\begin{solution}
设 \(M\) Noether、\(f\) 满。取 \(n\) 使 \(\ker f^n=\ker f^{n+1}\)。若 \(x\in\ker f\)，由 \(f^n\) 满可取 \(y\) 使 \(f^n(y)=x\)。于是 \(y\in\ker f^{n+1}=\ker f^n\)，故 \(x=0\)，\(f\) 单且满。

设 \(M\) Artin、\(f\) 单。取 \(n\) 使 \(\operatorname{im}f^n=\operatorname{im}f^{n+1}\)。对 \(x\in M\)，可取 \(y\) 使 \(f^n(x)=f^{n+1}(y)\)。因 \(f^n\) 单，\(x=f(y)\)，故 \(f\) 满。整数模上的乘二单而不满，给出第一个反向失败；Artin 模 \(C_{p^\infty}\) 上的乘 \(p\) 满而不单，因为每个分式可以再除以 \(p\)，其核为非零的 \(C_p\)。
\end{solution}

\begin{exercise}\label{b2:ex:12-left-right-length}
设域扩张 \(k\subseteq K\) 的次数为有限整数 \(d\)，\(r\ge1\) 为整数。取由三元组 \((a,\mathbf u,b)\in K\times K^r\times k\) 组成的环 \(R\)，加法逐项进行，乘法为
\[
(a,\mathbf u,b)(a',\mathbf u',b')
=(aa',a\mathbf u'+\mathbf u b',bb').
\]
单位元为 \((1,\mathbf0,1)\)。分别计算左右正则模的长度，并比较 \(k=\mathbb R,K=\mathbb C,r=2\) 的情形。
\end{exercise}
\begin{solution}
这个乘法来自 \(K^r\) 的左 \(K\)、右 \(k\) 标量作用；两侧作用相容，展开三元组乘法即可验证结合律。令 \(e_1=(1,\mathbf0,0)\)、\(e_2=(0,\mathbf0,1)\)，则它们正交且和为一。由此左右正则模分别分解为 \(Re_1\oplus Re_2\) 与 \(e_1R\oplus e_2R\)。

记 \(J=\{\,(0,\mathbf u,0):\mathbf u\in K^r\,\}\)。左乘在 \(J\) 上只作用为 \(\mathbf u\mapsto a\mathbf u\)，所以其子模恰为 \(K^r\) 的 \(K\) 子空间，\(\ell({}_RJ)=r\)。\(Re_1\) 是由左上角作用的简单模 \(K\)，\(Re_2/J\) 是由右下角作用的简单模 \(k\)。长度可加性给出左正则长度 \(1+r+1=r+2\)。

右乘在 \(J\) 上只作用为 \(\mathbf u\mapsto\mathbf u b\)，其子模恰为 \(K^r\) 的 \(k\) 子空间，故 \(\ell(J_R)=rd\)。\(e_2R\) 是简单右模 \(k\)，\(e_1R/J\) 是简单右模 \(K\)，于是右正则长度为 \(rd+2\)。取实复扩张及 \(r=2\)，得到左右长度分别为四和六。相比\cref{b2:prop:triangular-left-right-length}，扩大中间的双模后，这一部分贡献的左侧简单层数为 \(r\)，右侧为 \(rd\)。
\end{solution}

\begin{exercise}\label{b2:ex:12-cyclic-length}
求整数模 \(M=\mathbb Z/360\mathbb Z\) 的长度及各简单因子的重数，并计算乘十二自同态的核与像的长度。
\end{exercise}
\begin{solution}
\(360=2^3\cdot3^2\cdot5\)。互素循环分解给出三种主分量，分别有三、二、一层简单因子，所以 \(\ell(M)=6\)，\(\mathbb Z/2,\mathbb Z/3,\mathbb Z/5\) 的重数分别为三、二、一。乘十二的核由 \(\bar{30}\) 生成，阶为十二；像由 \(\bar{12}\) 生成，阶为三十。因此两者分别同构于 \(\mathbb Z/12\mathbb Z\)、\(\mathbb Z/30\mathbb Z\)，长度均为三。长度可加性得到 \(6=3+3\)，但这两个长度三的模具有不同的简单因子重数。
\end{solution}

\begin{exercise}\label{b2:ex:12-truncated-polynomial-length}
设 \(k\) 为域，\(A=k[x,y]/(x,y)^3\)，把 \(A\) 看成自身的幺性左模。写出一条合成列，求其长度与合成因子。计算乘 \(x\) 的核和像，并核验长度可加性。
\end{exercise}
\begin{solution}
以 \(1,x,y,x^2,xy,y^2\) 的像为基，记 \(\mathfrak m=(x,y)\)。一条合成列是
\[
0\subsetneq kx^2\subsetneq kx^2+kxy\subsetneq\mathfrak m^2
\subsetneq\mathfrak m^2+kx\subsetneq\mathfrak m\subsetneq A.
\]
各项都为理想：二次部分被 \(x,y\) 零化，而乘一次部分所得都在 \(\mathfrak m^2\)。每个相邻商一维，且 \(x,y\) 作用为零，故都是简单模 \(A/\mathfrak m\cong k\)。于是长度为六。

乘 \(x\) 将 \(1,x,y\) 分别送到 \(x,x^2,xy\)，将三个二次基元素送到零。因此其像以 \(x,x^2,xy\) 为基，核为 \(\mathfrak m^2\)，两者均为三维。它们的合成因子都只能为上述 \(k\)，从列中的相应细分可得长度均为三，因此 \(6=3+3\)。
\end{solution}

\begin{exercise}\label{b2:ex:12-matrix-module-length}
设 \(k\) 为域，\(n\ge1,r\ge0\) 为整数。令 \(A=M_n(k)\)，以左矩阵乘法作用在 \(M_{n\times r}(k)\) 上。求该模的长度与一条合成列，说明长度与 \(k\) 维数的关系。
\end{exercise}
\begin{solution}
按列分解为 \(r\) 个 \(k^n\) 的直和，每个列空间都是简单左 \(A\) 模：从非零列向量的一个非零坐标出发，矩阵单位及标量可得到全部标准列向量。逐次加入前一、前二、直到全部 \(r\) 列，得到一条有 \(r\) 个简单商层的合成列。因此长度为 \(r\)，而 \(k\) 维数为 \(nr\)。只在 \(n=1\) 或零模情形，两者数值相同。
\end{solution}

\begin{exercise}\label{b2:ex:12-fitting-integers}
在整数模 \(M=\mathbb Z/600\mathbb Z\) 上令 \(f\) 为乘六。求最小的正整数 \(n\)，使 \(M=\ker f^n\oplus\operatorname{im}f^n\)。写出两个分量的投影，求 \(f\) 在可逆分量上的逆，并计算两分量的长度。
\end{exercise}
\begin{solution}
分解 \(600=24\cdot25=2^3\cdot3\cdot5^2\)。乘六在二、三主分量上幂零，在五主分量上可逆。因 \(\gcd(6^3,600)=24\)，有
\[
\ker f^3=\langle\bar{25}\rangle\cong\mathbb Z/24\mathbb Z,
\qquad
\operatorname{im}f^3=\langle\bar{24}\rangle\cong\mathbb Z/25\mathbb Z.
\]
它们交为零，且 \(25-24=1\) 保证其和为 \(M\)。当 \(n=1,2\) 时，非零元 \(\bar{300}\) 都在 \(\ker f^n\) 中；它又是六与十二的倍数，而这两数分别是 \(\gcd(6,600)\)、\(\gcd(6^2,600)\)，所以也在相应的像中。因此最小正整数为三。

乘二十五在阶二十四的分量上为恒等，在阶二十五的分量上为零，故为核分量投影；其余投影为乘 \(1-25\)，即乘五百七十六。它们是和为恒等的正交幂等元。可逆分量上乘二十一是 \(f\) 的逆，因为 \(6\cdot21\equiv1\pmod{25}\)。两分量的长度分别为 \(3+1=4\) 与二，合为 \(\ell(M)=6\)。
\end{solution}

\begin{exercise}\label{b2:ex:12-fitting-polynomial}
设 \(k\) 为域，\(p,q\in k[t]\) 为互不相伴的首一不可约多项式，\(a,b\ge1\) 为整数。在
\[
M=k[t]/(p^a)\oplus k[t]/(q^b)
\]
上令 \(f\) 为乘 \(p\)。求其 Fitting 分解，并求使分解成立的最小正整数。
\end{exercise}
\begin{solution}
第一分量上 \(f^a=0\)，而 \(f^{a-1}\ne0\)；第二分量上 \(p\) 与 \(q^b\) 互素，Bézout 等式给出乘 \(p\) 的逆。因此 \(n=a\) 时，核为第一分量，像为第二分量，得到分解。

若 \(1\le n<a\)，第一分量的核为 \((p^{a-n})/(p^a)\)，像为 \((p^n)/(p^a)\)，两者交含非零的 \(p^{a-1}+(p^a)\)，因为 \(a-1\ge n,a-n\)。所以此时不能为直和，最小整数正是 \(a\)。第二分量的指数 \(b\) 不影响这个最小值，因为 \(f\) 在其上从一开始就可逆。
\end{solution}

\begin{exercise}\label{b2:ex:12-local-not-finite-length}
设 \(k\) 为域，\(R=k[\![t]\!]\) 为形式幂级数环。证明 \(a(t)=\sum_{n\ge0}a_nt^n\) 是单位当且仅当 \(a_0\ne0\)，并在单位情形递推求出逆级数的系数。证明 \(R\) 局部，其左正则模有限生成且不可分解，却不有限长度；给出不幂零的非单位。
\end{exercise}
\begin{solution}
若 \(a(t)b(t)=1\)，常数项满足 \(a_0b_0=1\)，所以 \(a_0\ne0\)。反之，令
\[
b_0=a_0^{-1},\qquad
b_n=-a_0^{-1}\sum_{i=1}^n a_i b_{n-i}\quad(n\ge1).
\]
每一项只使用已确定的有限个系数，使乘积的常数项为一、其他系数为零，故得到形式幂级数逆。这里不涉及级数的解析收敛。

非单位恰为常数项为零的级数，即真理想 \(tR\)，所以 \(R\) 局部。左正则模由一生成；其自同态环由右乘法识别为 \(R^{\mathrm{op}}=R\)，因而局部，由\cref{b2:prop:local-ring-criterion}中幂等元判据及\cref{b2:prop:idempotent-direct-sum}得不可分解。另一方面，\(R\supsetneq tR\supsetneq t^2R\supsetneq\cdots\) 严格：\(t^n\) 的 \(n\) 次系数非零，而 \(t^{n+1}R\) 中该系数全为零。因此它不 Artin，也不有限长度。非单位 \(t\) 的每次幂都非零，故不幂零。
\end{solution}

\begin{exercise}\label{b2:ex:12-jordan-centralizer}
设 \(k\) 为域，\(n\ge1\)，\(J=J_n(0)\in M_n(k)\) 为单个上次对角线为一的幂零 Jordan 块。证明与 \(J\) 交换的全部矩阵恰为
\[
a_0I+a_1J+\cdots+a_{n-1}J^{n-1}.
\]
证明这个环局部，并确定其中的幂等元。
\end{exercise}
\begin{solution}
向量 \(e_n\) 在 \(J\) 下循环，因为 \(e_n,Je_n,\ldots,J^{n-1}e_n\) 是逆序标准基。代入映射 \(P(t)\mapsto P(J)e_n\) 因而给出 \(k[t]/(t^n)\cong k^n\)：\(J^n=0\)，而上述基保证次数小于 \(n\) 的非零多项式不在核中。特别地，\(\mu_J(t)=t^n\)，\(k^n\) 在作用代数 \(k[J]\) 下同构于其正则模。由\cref{b2:thm:regular-double-centralizer}，正则左作用的交换子来自右乘；\(k[J]\) 交换，右乘就是左乘，所以 \(\operatorname{Cent}(J)=k[J]\)。上述基也保证 \(I,J,\ldots,J^{n-1}\) 的线性组合表示唯一。

常数项 \(a_0\ne0\) 时，提出 \(a_0\) 后用幂零矩阵的有限几何和求逆；常数项为零时，矩阵为 \(J\) 的倍数，因而幂零、非可逆。非单位构成理想 \((J)\)，所以环局部。由\cref{b2:prop:local-ring-criterion}，幂等元只有零和单位矩阵。
\end{solution}

\begin{exercise}\label{b2:ex:12-nilpotent-sum}
设 \(k\) 为域。在 \(M_2(k)\) 中给出两个幂零矩阵，其和可逆。说明为什么这不违背有限长度不可分解模的局部自同态环结论。
\end{exercise}
\begin{solution}
取 \(U=E_{12}\)、\(V=E_{21}\)，有 \(U^2=V^2=0\)，但 \((U+V)^2=I\)，故其和可逆，任意特征下都成立。它们是 \(k^2\) 的自同态，而这个模可分解为两个一维子空间。定理要求有限长度与不可分解同时成立；在满足这两个条件的同一个模上，非可逆自同态才组成理想，幂零元素之和才由该理想性继续非可逆并因 Fitting 判据而幂零。
\end{solution}

\begin{exercise}\label{b2:ex:12-graph-projection}
设 \(R\) 为含幺环，\(M=A\oplus B\) 为幺性左 \(R\) 模的直和，\(u:A\rightarrow B\) 为 \(R\) 模同态。证明 \(u\) 的图
\(\Gamma_u=\Bigl\{\,\bigl(a,u(a)\bigr)\mid a\in A\,\Bigr\}\)
是以 \(B\) 为补模的直和因子，求对应投影，并证明它与投影到 \(A\) 的幂等元在 \(\operatorname{End}_R(M)\) 中共轭。
\end{exercise}
\begin{solution}
每个 \((a,b)\) 唯一写成 \(\bigl(a,u(a)\bigr)+\bigl(0,b-u(a)\bigr)\)，故 \(M=\Gamma_u\oplus B\)。对应投影为
\[
e_u=\begin{pmatrix}1&0\\u&0\end{pmatrix}.
\]
令 \(g=\begin{pmatrix}1&0\\u&1\end{pmatrix}\)，则 \(g^{-1}=\begin{pmatrix}1&0\\-u&1\end{pmatrix}\)，并且
\[
g\begin{pmatrix}1&0\\0&0\end{pmatrix}g^{-1}=e_u.
\]
所以不同图子模虽然作为原模中的子模位置不同，其投影通过这个明确自同构共轭，图本身也都同构于 \(A\)。
\end{solution}

\begin{exercise}\label{b2:ex:12-power-cancellation}
设 \(R\) 为含幺环，\(M,N\) 为有限长度幺性左 \(R\) 模，\(r\ge1\) 为整数。证明 \(M^{\oplus r}\cong N^{\oplus r}\) 蕴含 \(M\cong N\)。这里不能未经说明就逐次消去，因为两侧并未给定相同的直和因子。
\end{exercise}
\begin{solution}
先把 \(M,N\) 分解为有限个不可分解模。对每一种出现的同构类型，设其重数分别为 \(a,b\)。两边的 \(r\) 次直和中相应重数为 \(ra,rb\)，由\cref{b2:thm:krull-schmidt}必相等。正整数 \(r\) 可在整数等式中约去，所以 \(a=b\)。全部类型的重数都相同，逐项同构的直和便给出 \(M\cong N\)。这里约去的是整数重数，不是预先假定的模消去律。
\end{solution}

\begin{exercise}\label{b2:ex:12-finite-extension-descent}
设 \(L/k\) 为有限域扩张，\(A\) 为结合含幺 \(k\) 代数，\(M,N\) 为有限维幺性左 \(A\) 模。若 \(L\otimes_kM\cong L\otimes_kN\) 作为 \(L\otimes_kA\) 模，证明 \(M\cong N\) 作为 \(A\) 模。可对照有限扩张情形的证明方法\cite[Problem 3.8.4(i), pp. 52--53]{EtingofEtAl2011RepresentationTheory}。
\end{exercise}
\begin{solution}
令 \(d=[L:k]\)，选择 \(L\) 的 \(k\) 基 \(\lambda_1,\ldots,\lambda_d\)。限制为 \(A\) 模后，映射
\[
M^{\oplus d}\longrightarrow L\otimes_kM,\qquad
(m_1,\ldots,m_d)\longmapsto\sum_i\lambda_i\otimes m_i
\]
是同构：张量积的基展开给出唯一坐标，\(A\) 只作用在第二因子，所以映射 \(A\) 线性。对 \(N\) 同样成立，题设因而给出 \(M^{\oplus d}\cong N^{\oplus d}\)。两模有限维，其子模链由维数控制，故有限长度。分解为不可分解模后，Krull–Schmidt 定理使每一种类型的重数满足 \(d a=d b\)，故 \(a=b\)，得到 \(M\cong N\)。有限扩张条件用于把限制标量后的模写成有限个相同分量。
\end{solution}

\begin{exercise}\label{b2:ex:12-infinite-cancellation}
设 \(k\) 为域，\(k\) 向量空间 \(V\) 有可数无限基 \(e_0,e_1,\ldots\)。构造 \(k\oplus V\cong V\)，据此说明去掉有限长度以后，\(A\oplus C\cong B\oplus C\) 不再保证 \(A\cong B\)。
\end{exercise}
\begin{solution}
把左侧 \(k\) 的标准基向量送到 \(e_0\)，把 \(V\) 中的 \(e_i\) 送到 \(e_{i+1}\)，得到两组基之间的双射，线性延拓即为同构，逆映射将 \(e_0\) 送回首个分量，将其余基向量依次前移。取 \(A=k,B=0,C=V\)，则 \(A\oplus C\cong B\oplus C\)，但 \(k\not\cong0\)。无限公共分量能够吸收新增的一维空间，整数重数的有限消去论证不再适用。
\end{solution}

\begin{exercise}\label{b2:ex:12-nonprincipal-projective}
对\cref{b2:ex:krull-schmidt-nonunique}中的 \(R=\mathbb Z[\sqrt{-5}]\)、\(I=(2,1+\sqrt{-5})\)，证明 \(I\) 是有限生成投射模，却不是自由模。
\end{exercise}
\begin{solution}
\(I\) 由题述两个元素生成。正文给出的显式同构 \(I\oplus I\cong R^2\) 说明 \(I\) 是自由模的直和因子，由\cref{b2:thm:projective-characterizations}，它投射。若 \(I\) 自由，因为非零，其基非空。任意两个非零元素 \(a,b\in I\) 都满足 \(R\) 线性关系 \(b\cdot a+(-a)\cdot b=0\)。这里 \(a,b\) 首先是模元素，又因 \(I\subseteq R\) 而能作标量；系数 \(b,-a\) 均非零，所以二者不能同时属于一组基。基因而至多有一个元素。单元素基却意味着 \(I\) 为主理想，与正文已经证明的非主性矛盾。因此它不自由。
\end{solution}

\begin{exercise}\label{b2:ex:12-dual-numbers-decomposition}
设 \(k\) 为域，\(A=k[\varepsilon]/(\varepsilon^2)\)，\(M\) 为五维幺性左 \(A\) 模，乘 \(\varepsilon\) 的线性算子秩为二。确定 \(M\) 的不可分解直和分解及长度。
\end{exercise}
\begin{solution}
记算子为 \(T\)，则 \(T^2=0\)。由\cref{b2:thm:jordan-canonical-form}，其 Jordan 块只能为大小一或二的零特征值块，因为大小至少三的块平方不为零。每个二阶块贡献秩一，每个一阶块贡献秩零，所以恰有两个二阶块和一个一阶块。二阶块作为 \(A\) 模同构于正则模 \(A\)，一阶块同构于 \(A/(\varepsilon)\cong k\)，故
\[
M\cong A\oplus A\oplus k.
\]
\(A\) 的非零子模都包含 \(k\varepsilon\)，所以不可分解，\(k\) 简单也不可分解。长度分别为二、二、一，总长度为五，三个不可分解因子的个数则为三。
\end{solution}

\begin{exercise}\label{b2:ex:12-orthogonal-family-bound}
设 \(R\) 为含幺环，\(M\) 为长度 \(d>0\) 的幺性左 \(R\) 模。若 \(\operatorname{End}_R(M)\) 中有 \(d\) 个两两正交的非零幂等元 \(e_1,\ldots,e_d\)，证明它们的和自动为恒等，每个 \(e_iM\) 都简单。反之，若 \(M\) 是简单模的有限直和，证明正交非零幂等元的个数可以达到 \(d\)。比较整数模 \(\mathbb Z/p^2\mathbb Z\) 与 \((\mathbb Z/p\mathbb Z)^2\) 能达到的最大个数，其中 \(p\) 为素数。
\end{exercise}
\begin{solution}
令 \(e=\sum_i e_i\)。由正交性，\(e^2=e\)，且 \(eM=\bigoplus_i e_iM\)。长度可加性给出
\[
d\le\sum_{i=1}^d\ell(e_iM)=\ell(eM)\le\ell(M)=d.
\]
所以每项长度为一，因而简单；\(\ell(M/eM)=0\) 则迫使 \(eM=M\)。幂等元在自身像上为恒等，故 \(e=1_M\)。

若 \(M\) 为 \(s\) 个简单模的直和，其长度为 \(s=d\)，各分量投影便给出 \(d\) 个正交非零幂等元。对 \(\mathbb Z/p^2\mathbb Z\)，任何非零子模都包含其唯一的阶 \(p\) 子模；两非零直和因子不可能交于零，故它不可分解，自同态环只有零和恒等两个幂等元，最大个数为一。\((\mathbb Z/p\mathbb Z)^2\) 的两个坐标投影达到二，且\cref{b2:prop:orthogonal-idempotent-decomposition}的长度界禁止更多。这两个模的简单商层相同，但实际可分离的投影数不同。
\end{solution}
